\chapter{Πειράματα και Μετρήσεις}
\label{chap:experiments}

Το παρόν κεφάλαιο παρουσιάζει το πειραματικό σκέλος της εργασίας: το κύριο σώμα μετρήσεων επιδόσεων (\eng{benchmark}), τα τεκμήρια ακριβούς αναζήτησης (\eng{exact-search evidence}) των εγγενών και των ενσωματωμένων επιλυτών, την εγγενή απόδειξη βελτιστότητας του superflip και την πλήρη εξαντλητική μελέτη του κύβου 2×2×2 (\eng{Pocket Cube}). Όλοι οι πίνακες του κεφαλαίου παράγονται αυτόματα από αποθηκευμένα αρχεία αποτελεσμάτων και κάθε αριθμός ανάγεται σε συγκεκριμένο αποθηκευμένο τεκμήριο (\eng{artifact}).

\section{Πειραματική διάταξη}
\label{sec:experiments-setup}

Όλες οι μετρήσεις εκτελέστηκαν σε ένα συμβατικό μηχάνημα με επεξεργαστή Apple M3 (8 πύρηνες) και 16~GiB μνήμης, με λειτουργικό σύστημα macOS. Οι εγγενείς μηχανές C/C++ μεταγλωττίζονται τοπικά, ενώ το στρώμα Python απαιτεί Python~3.10 ή νεότερη έκδοση. Οι περισσότερες εγγενείς εκτελέσεις χρησιμοποιούν 8 νήματα· όπου μια μέτρηση χρησιμοποιεί διαφορετικό πλήθος νημάτων, αυτό δηλώνεται ρητά στο κείμενο και στο αντίστοιχο τεκμήριο.

Οι χρονομετρήσεις πραγματοποιήθηκαν σε μηχάνημα υπό φορτίο γενικής χρήσης και όχι σε απομονωμένο περιβάλλον μετρήσεων. Για τον λόγο αυτό, οι χρονικά ευαίσθητες συγκρίσεις ακολουθούν δίκαιη λογιστική με κοινή προθέρμανση (\eng{warm-up}) και σταθερή, καταγεγραμμένη σειρά σκελών (βλ.~Ενότητα~\ref{sec:experiments-h48-fair}). Επιπλέον, επιλεγμένες μετρήσεις επαναλήφθηκαν σε συνθήκες χαμηλού φορτίου και αποθηκεύονται ως χωριστά τεκμήρια με κατάληξη \code{lowload}. Το κείμενο δίνει συστηματικά μεγαλύτερη έμφαση στις καταστάσεις αποτελέσματος και στις εγγυήσεις παρά σε μικρές διαφορές χιλιοστών του δευτερολέπτου.

Κάθε γραμμή αποτελεσμάτων φέρει μια κατάσταση αποτελέσματος (\eng{status}): \code{exact} σημαίνει αποδεδειγμένα βέλτιστη και ανεξάρτητα επαληθευμένη λύση, \code{non_exact} σημαίνει επαληθευμένη λύση χωρίς απόδειξη βελτιστότητας, \code{lower_bound} σημαίνει αποδεδειγμένο κάτω φράγμα απόστασης χωρίς λύση, και \code{timeout} σημαίνει μη ολοκλήρωση εντός του χρονικού ορίου χωρίς αποτέλεσμα. Το πλήρες σχήμα των αρχείων αποτελεσμάτων ορίζεται στο Παράρτημα~\ref{app:schema}.

Τα χρονικά όρια διαφέρουν ανά οικογένεια μετρήσεων και καταγράφονται στα αντίστοιχα τεκμήρια:
\begin{itemize}
  \item Κύριο σώμα μετρήσεων: ο οριοθετημένος Korf/IDA* σε Python εκτελείται με όριο βάθους 10· η αποθηκευμένη εκτέλεση δεν παρήγαγε καμία γραμμή \code{timeout}, και οι βαθύτερες μη ολοκληρώσεις καταγράφονται ως \code{lower_bound}.
  \item Μετρήσεις πίεσης H48: όριο 5~s ανά περίπτωση και σκληρό όριο 20~s για τον πίνακα \code{h48h0}· όριο 10~s και σκληρό όριο 60~s για τον πίνακα \code{h48h7}.
  \item Πιστοποιήσεις του μαντείου (\eng{oracle}): όριο 90~s για την ελεγχόμενη λειτουργία, 180~s για την έμπιστη λειτουργία με προφόρτωση, 300~s για την έμπιστη λειτουργία χωρίς προφόρτωση και 240~s για την πιστοποίηση με παραμένουσα διεργασία.
  \item Συγκριτικές μικρομετρήσεις: όριο 120~s για τη σύγκριση ελεγχόμενης/έμπιστης λειτουργίας, 180~s για τη λειτουργία δέσμης, 30~s για την παραμένουσα διεργασία και 60~s για τις διεπαφές γραμμής εντολών, ροής και προγραμματιστικής χρήσης.
  \item Εξωτερική διαδρομή Nissy: όριο 60~s ανά περίπτωση με σκληρό όριο 120~s στο σώμα πίεσης· οι στοχευμένες δοκιμές superflip χρησιμοποιούν όρια 120~s με 2 νήματα και 240~s με 8 νήματα.
\end{itemize}

Τα κυριότερα αρχεία τεκμηρίων του κεφαλαίου είναι τα εξής:
\begin{itemize}
  \item Πρωτογενείς γραμμές μετρήσεων: \path{results/raw/benchmarks_seed_2026_thesis.jsonl}.
  \item Επεξεργασμένη σύνοψη: \path{results/processed/summary_seed_2026_thesis.json}.
  \item Τεκμήρια εγγενούς βέλτιστης αναζήτησης: \path{results/processed/optimal_3x3_seed_2026_thesis.json}.
  \item Τεκμήρια πίεσης H48: \path{results/processed/optimal_3x3_seed_2026_stress_h48h0.json} και \path{results/processed/optimal_3x3_seed_2026_stress_h48h7_oracle.json}.
  \item Δίκαια συγκριτικά τεκμήρια: \path{results/processed/h48_trusted_table_speedup_seed_2026_thesis_h48h7_fair.json}, \path{results/processed/h48_batch_overhead_seed_2026_thesis_trusted_fair.json} και \path{results/processed/h48_resident_oracle_seed_2026_thesis_h48h7_trusted_fair.json}.
\end{itemize}
Οι πλήρεις εντολές παραγωγής και αναπαραγωγής όλων των μετρήσεων παρατίθενται στο Παράρτημα~\ref{app:repro} και στο Παράρτημα~\ref{app:cli}, ενώ αναλυτικά παραδείγματα χρήσης δίνονται στο Κεφάλαιο~\ref{chap:usage}.

\section{Μεθοδολογία}
\label{sec:experiments-methodology}

Το κύριο σώμα μετρήσεων χρησιμοποιεί σπόρο (\eng{seed}) 2026, τη λυμένη κατάσταση, ρηχές ντετερμινιστικές διαταράξεις (\eng{scrambles}) βάθους 1 έως 8 και ντετερμινιστικές τυχαίες διαταράξεις ονομαστικού βάθους 5, 10, 15 και 20. Τα πρωτογενή αποτελέσματα αποθηκεύονται σε μορφή JSONL και η επεξεργασμένη σύνοψη σε JSON. Οι πίνακες που ακολουθούν παράγονται από δέσμες ενεργειών (\eng{scripts}) και δεν γράφτηκαν χειροκίνητα. Το προφίλ \code{quick} χρησιμοποιείται μόνο για ανάπτυξη και γρήγορους ελέγχους και δεν τεκμηριώνει κανέναν ισχυρισμό του κεφαλαίου.

Η επιλογή σταθερού σπόρου επιτρέπει την επανάληψη του ίδιου πειράματος χωρίς χειροκίνητη αποθήκευση κάθε τυχαίας απόφασης. Οι ρηχές περιπτώσεις έχουν γνωστό ονομαστικό βάθος διατάραξης, όμως το ονομαστικό βάθος δεν ταυτίζεται πάντα με την ακριβή απόσταση: μια ακολουθία διατάραξης μπορεί να εμφανίζει ακυρώσεις ή να επιδέχεται συντομότερη λύση. Για τον λόγο αυτό, όπου είναι εφικτό, η ακριβής απόσταση ελέγχεται με BFS και αποθηκεύεται στο πεδίο \code{known_exact_distance}.

Το κύριο σώμα μετρήσεων δεν επιδιώκει να εξαντλήσει τον χώρο καταστάσεων του κύβου. Στόχος του είναι ένα μικρό, ελεγχόμενο και αναπαραγώγιμο σύνολο περιπτώσεων που αναδεικνύει τη συμπεριφορά των επιλυτών σε λυμένη, ρηχή και βαθύτερη κατάσταση, ώστε κάθε γραμμή να μπορεί να ελεγχθεί και να συζητηθεί χωριστά. Για στατιστικά ισχυρό συμπέρασμα γενικής απόδοσης θα απαιτούνταν πολύ μεγαλύτερο σώμα, περισσότερες επαναλήψεις και πιθανώς διαφορετική υπολογιστική υποδομή.

Οι μετρήσεις χρόνου εκτέλεσης πρέπει να διαβάζονται ως τοπικές ενδείξεις και όχι ως καθολική κατάταξη αλγορίθμων. Επηρεάζονται από το υλικό, τον διερμηνευτή της Python, τη θέρμανση της κρυφής μνήμης, τα προαιρετικά πακέτα και τα καθορισμένα χρονικά όρια. Για τον λόγο αυτό το κείμενο δίνει μεγαλύτερη βαρύτητα στις καταστάσεις αποτελέσματος, στην επαλήθευση και στη διάκριση \code{exact}/\code{non_exact} παρά σε μικρές χρονικές διαφορές.

\section{Σύνολα περιπτώσεων}
\label{sec:experiments-cases}

Το σύνολο A περιέχει τη λυμένη κατάσταση, το σύνολο B τις ρηχές ντετερμινιστικές διαταράξεις βάθους 1 έως 8 και το σύνολο C τις ντετερμινιστικές τυχαίες περιπτώσεις ονομαστικού βάθους 5 έως 20. Το σύνολο C περιλαμβάνει δηλαδή και μία σχετικά ρηχή τυχαία περίπτωση βάθους 5, ενώ οι περιπτώσεις ονομαστικού βάθους 15 και 20 λειτουργούν ως πρακτικές δοκιμές πίεσης. Η διάκριση αυτή επιτρέπει να διαβάζεται χωριστά η ορθότητα σε ελεγχόμενα βάθη και η πρακτική συμπεριφορά σε βαθύτερες διαταράξεις. Ο Πίνακας~\ref{tab:benchmark-datasets} συνοψίζει τα τρία σύνολα.

Η λυμένη κατάσταση αποτελεί στοιχειώδη έλεγχο ορθότητας: ένας επιλυτής που δεν χειρίζεται σωστά την ταυτότητα δεν μπορεί να θεωρηθεί αξιόπιστος σε δυσκολότερες περιπτώσεις. Οι ρηχές περιπτώσεις του συνόλου B ελέγχουν τη σχέση επιλυτή και επαληθευτή σε καταστάσεις όπου η απόσταση αποδεικνύεται με μικρή αναζήτηση. Στις βαθύτερες περιπτώσεις του συνόλου C αναμένεται ότι οι ακριβείς επιλυτές μπορεί να φθάσουν στα όριά τους, ενώ οι πρακτικοί επιλυτές μπορεί να βρουν λύσεις χωρίς απόδειξη βελτιστότητας.

\begin{table}[htbp]\centering
\begin{tabular}{lrrr}
\hline
Dataset & Cases & Min depth & Max depth \\
\hline
A & 1 & 0 & 0 \\
B & 8 & 1 & 8 \\
C & 4 & 5 & 20 \\
\hline
\end{tabular}

\caption[Σύνοψη των συνόλων περιπτώσεων του κύριου σώματος μετρήσεων]{Σύνοψη των συνόλων περιπτώσεων A, B και C του κύριου σώματος μετρήσεων: πλήθος περιπτώσεων και εύρος ονομαστικού βάθους διατάραξης. Πηγή: \path{results/raw/benchmarks_seed_2026_thesis.jsonl}.}
\label{tab:benchmark-datasets}
\end{table}

Ο πλήρης κατάλογος των περιπτώσεων καταγράφεται στον Πίνακα~\ref{tab:benchmark-catalog}, ώστε κάθε γραμμή μετρήσεων να μπορεί να αναπαραχθεί από τον ίδιο σπόρο και την ίδια ακολουθία κινήσεων.

\begin{table}[htbp]\centering
\begin{tabular}{p{0.22\textwidth}lrp{0.42\textwidth}}
\hline
Case & Dataset & Depth & Scramble \\
\hline
solved & A & 0 & -- \\
shallow\_1 & B & 1 & U' \\
shallow\_2 & B & 2 & R' U' \\
shallow\_3 & B & 3 & U' B' L' \\
shallow\_4 & B & 4 & D R2 B' D2 \\
shallow\_5 & B & 5 & L' B' L' F' R \\
shallow\_6 & B & 6 & D2 F2 L' B' F' B \\
shallow\_7 & B & 7 & U F2 D' U2 B' R F \\
shallow\_8 & B & 8 & B2 D F' R D L' B L \\
random\_0\_5 & C & 5 & D2 R' B2 U R' \\
random\_1\_10 & C & 10 & U' L' U B2 F L B2 R' L2 R' \\
random\_2\_15 & C & 15 & R' D R2 D' B' D L2 F' B2 F2 L' F U2 R2 L \\
random\_3\_20 & C & 20 & L D2 U2 B' U L' B' U R F2 U2 L2 B2 R F L D R L D' \\
\hline
\end{tabular}

\caption[Κατάλογος περιπτώσεων του κύριου σώματος μετρήσεων]{Κατάλογος των 13 περιπτώσεων του κύριου σώματος μετρήσεων, με το σύνολο, το ονομαστικό βάθος και την πλήρη ακολουθία διατάραξης κάθε περίπτωσης. Πηγή: \path{results/raw/benchmarks_seed_2026_thesis.jsonl}.}
\label{tab:benchmark-catalog}
\end{table}

\section{Έλεγχος από άκρο σε άκρο για τον κύβο 3×3×3}
\label{sec:experiments-e2e}

Πέρα από το κύριο σώμα μετρήσεων, εκτελείται ρητός έλεγχος από άκρο σε άκρο (\eng{end-to-end}) για πραγματικές εισόδους 3×3×3: λυμένη κατάσταση, μικρή ακολουθία κινήσεων, βαθύτερη ακολουθία 20 κινήσεων και η ίδια βαθύτερη κατάσταση ως συμβολοσειρά ψηφίδων (\eng{facelets}). Για κάθε γραμμή εκτελούνται έλεγχος εγκυρότητας, αναγνώριση απόστασης ή κάτω φράγματος, αυτόματη πρακτική επίλυση και ανεξάρτητη επαλήθευση της λύσης. Η αυτόματη πρακτική επίλυση δοκιμάζει πρώτα την εγγενή οριοθετημένη διαδρομή δύο φάσεων Kociemba, με προεπιλεγμένα όρια βάθους ίσα με τις φασικές διαμέτρους (12 για τη φάση 1 και 18 για τη φάση 2). Έτσι και οι βαθύτερες γραμμές των 20 κινήσεων λύνονται από την εγγενή διαδρομή και όχι από τον εξωτερικό προσαρμογέα (\eng{adapter}), ενώ για τη βαθιά κατάσταση η αναγνώριση απόστασης επιστρέφει τεκμηριωμένο κάτω φράγμα. Ο έλεγχος αυτός δεν μετατρέπει τα \code{non_exact} αποτελέσματα σε βέλτιστα· τεκμηριώνει όμως ότι η διαδρομή 3×3×3 του συστήματος λειτουργεί από άκρο σε άκρο τόσο για είσοδο ακολουθίας όσο και για είσοδο ψηφίδων. Τα αποτελέσματα συνοψίζονται στον Πίνακα~\ref{tab:e2e-status}· η εντολή εκτέλεσης παρατίθεται στο Παράρτημα~\ref{app:cli}.

\begin{table}[htbp]\centering
{\small
\begin{tabular}{L{0.20\textwidth}L{0.12\textwidth}L{0.18\textwidth}L{0.14\textwidth}L{0.12\textwidth}r}
\hline
Case & Input & Distance & Auto path & Status & Length \\
\hline
solved & solved & exact\_distance & Koc. native & non\_exact & 0 \\
shallow\_sequence & sequence & exact\_distance & Koc. native & non\_exact & 3 \\
deep\_sequence\_20 & sequence & lower\_bound & Koc. native & non\_exact & 21 \\
deep\_facelets\_20 & facelets & lower\_bound & Koc. native & non\_exact & 21 \\
\hline
\end{tabular}
}

\caption[Αποτελέσματα ελέγχου από άκρο σε άκρο για τον κύβο 3×3×3]{Αποτελέσματα του ελέγχου από άκρο σε άκρο για πραγματικές εισόδους 3×3×3: είδος εισόδου, είδος αναγνώρισης απόστασης, διαδρομή αυτόματης επίλυσης, κατάσταση αποτελέσματος και μήκος λύσης. Πηγή: τεκμήριο του \code{scripts/run_3x3_end_to_end.py} για το προφίλ \code{thesis}.}
\label{tab:e2e-status}
\end{table}

\section{Αναγνώριση απόστασης}
\label{sec:experiments-distance}

Η εντολή αναγνώρισης απόστασης εκτελεί πρώτα ρηχή εξαντλητική αναζήτηση όπου αυτό είναι εφικτό· διαφορετικά επιστρέφει τεκμηριωμένο κάτω φράγμα ή άγνωστο αποτέλεσμα λόγω ορίου. Ο Πίνακας~\ref{tab:distance-recognition} συνοψίζει το είδος απόκρισης στις μοναδικές περιπτώσεις του σώματος μετρήσεων, όχι ανά επιλυτή.

Η καταμέτρηση ανά μοναδική περίπτωση είναι σημαντική, επειδή η απόσταση είναι ιδιότητα της κατάστασης και όχι του επιλυτή. Αν τέσσερις επιλυτές εκτελούνται στην ίδια κατάσταση, το ίδιο ερώτημα αναγνώρισης απόστασης δεν πρέπει να μετρηθεί τέσσερις φορές. Το σώμα μετρήσεων κρατά γραμμές ανά επιλυτή για λόγους σύγκρισης, αλλά ο πίνακας αναγνώρισης απόστασης συνοψίζει τις αρχικές καταστάσεις, ώστε να αποφεύγεται η τεχνητή διόγκωση των πληθών.

Όταν η αναγνώριση επιστρέφει κάτω φράγμα, το αποτέλεσμα είναι χρήσιμο αλλά περιορισμένο. Ένα κάτω φράγμα 5 σημαίνει ότι δεν μπορεί να υπάρχει λύση μήκους μικρότερου από 5 σύμφωνα με τη συγκεκριμένη παραδεκτή εκτίμηση· δεν σημαίνει ότι υπάρχει λύση μήκους 5. Στο κείμενο και στα αρχεία αποτελεσμάτων η τιμή αυτή παραμένει διακριτή από την ακριβή απόσταση.

\begin{table}[htbp]\centering
\begin{tabular}{p{0.24\textwidth}p{0.36\textwidth}r}
\hline
Distance kind & Method & Cases \\
\hline
exact\_distance & bfs\_depth\_5 & 7 \\
exact\_distance & ida\_star\_depth\_10 & 4 \\
lower\_bound & combined\_table\_lower\_bound & 2 \\
\hline
\end{tabular}

\caption[Σύνοψη αναγνώρισης απόστασης ανά μοναδική περίπτωση]{Σύνοψη της αναγνώρισης απόστασης ανά μοναδική περίπτωση του κύριου σώματος μετρήσεων: είδος αποτελέσματος, μέθοδος και πλήθος περιπτώσεων. Πηγή: \path{results/processed/summary_seed_2026_thesis.json}.}
\label{tab:distance-recognition}
\end{table}

\section{Αξιολόγηση του εγγενούς βέλτιστου επιλυτή}
\label{sec:experiments-native-optimal}

Η κατασκευή του γωνιακού πίνακα προτύπων (\eng{pattern database}, PDB) και των οκτώ εξακμικών πινάκων προτύπων, μαζί με τα μεταδεδομένα και τους ελέγχους κάλυψης, περιγράφεται στο Κεφάλαιο~\ref{chap:implementation}. Η παρούσα ενότητα αξιολογεί τη συμπεριφορά της στοίβας ευρετικών και του εγγενούς βέλτιστου επιλυτή στο αποθηκευμένο σώμα τεκμηρίων.

Ο Πίνακας~\ref{tab:heuristic-comparison} παρουσιάζει τη στοίβα κάτω φραγμάτων που χρησιμοποιεί η IDA*. Στο ντετερμινιστικό δείγμα υπάρχουν 7 ρηχές γραμμές με ακριβή απόσταση από BFS και 6 βαθύτερες γραμμές σύγκρισης. Οι ρηχές γραμμές περνούν τον έλεγχο παραδεκτότητας και το τελικό συνδυασμένο κάτω φράγμα ισούται πάντα με το μέγιστο των προεπιλεγμένων συνιστωσών του. Η χρήση του μεγίστου διατηρεί την παραδεκτότητα, αν και οι προβολές των πινάκων επικαλύπτονται.

\begin{table}[htbp]\centering
{\scriptsize
\setlength{\tabcolsep}{3pt}
\begin{tabular}{lrrrrrrrrr}
\hline
Corpus & Cases & Exact & Mis. & Co. & Cor. & Ed. & Comb. & CPDB & Adm. \\
\hline
deterministic\_sample & 6 & 0 & 3 & 3.833 & 6.833 & 7.5 & 7.833 & 3.333 & -- \\
shallow\_exact & 7 & 7 & 1.714 & 1.571 & 2.571 & 2.714 & 2.714 & 1.429 & True \\
\hline
\end{tabular}
}

\caption[Σύγκριση συνιστωσών κάτω φράγματος της IDA*]{Σύγκριση των συνιστωσών κάτω φράγματος της IDA* (εσφαλμένα τοποθετημένα κυβίδια, γωνιακές και ακμικές προβολές, γωνιακή PDB, συνδυασμένο φράγμα) σε ρηχό και ντετερμινιστικό δείγμα, με έλεγχο παραδεκτότητας στις ρηχές γραμμές. Πηγή: τεκμήριο σύγκρισης ευρετικών του προφίλ \code{thesis}.}
\label{tab:heuristic-comparison}
\end{table}

Για να μη συγχέεται η πρακτική επίλυση με την ακριβή απόδειξη, η ακριβής αξιολόγηση εκτελείται ως χωριστή δέσμη μετρήσεων (βλ.~Παράρτημα~\ref{app:repro}). Ο εγγενής επιλυτής πλήρους κύβου με τον γωνιακό και τους οκτώ εξακμικούς πίνακες προτύπων χαρακτηρίζει μια γραμμή \code{exact} μόνο όταν η αναζήτηση ολοκληρώνεται και η λύση επαληθεύεται ξανά στον πλήρη κύβο. Όπως δείχνει ο Πίνακας~\ref{tab:optimal-status}, η διαδρομή αυτή αποδεικνύει ακριβείς λύσεις για τη λυμένη κατάσταση, τη ρηχή ακολουθία, την τυχαία περίπτωση ονομαστικού βάθους 10 (απόσταση 9) και την τυχαία περίπτωση ονομαστικού βάθους 15 (απόσταση 14).

\begin{table}[htbp]\centering
{\small
\begin{tabular}{lrrrrr}
\hline
Case & Depth & Initial LB & Status & Length & Seconds \\
\hline
solved & 0 & 0 & exact & 0 & 1.78449 \\
shallow\_sequence & 3 & 3 & exact & 3 & 1.57298 \\
random\_1\_10 & 10 & 7 & exact & 9 & 1.48651 \\
random\_2\_15 & 15 & 9 & exact & 14 & 20.0363 \\
\hline
\end{tabular}
}

\caption[Αποτελέσματα του εγγενούς βέλτιστου επιλυτή]{Αποτελέσματα του εγγενούς βέλτιστου επιλυτή με τον γωνιακό και τους οκτώ εξακμικούς πίνακες προτύπων: αρχικό κάτω φράγμα, κατάσταση αποτελέσματος, μήκος λύσης και χρόνος. Πηγή: \path{results/processed/optimal_3x3_seed_2026_thesis.json}.}
\label{tab:optimal-status}
\end{table}

\section{Τεκμήρια ακριβούς αναζήτησης με το H48}
\label{sec:experiments-h48}

Για τις βαθύτερες περιπτώσεις, όπου ο απλός εγγενής IDA* φθάνει σε πρακτικά όρια, η εργασία ενσωματώνει τη διαδρομή H48 του nissy-core~\cite{trontoNissyCoreH48} ως ενσωματωμένο (\eng{vendored}) εγγενή κώδικα με δικό της περίβλημα (\eng{wrapper}). Τα ευρήματα της ενότητας συνοψίζονται σε τρεις παρατηρήσεις. Πρώτον, η διαδρομή H48 ολοκληρώνει ακριβείς, ανεξάρτητα επαληθευμένες λύσεις για όλο το αποθηκευμένο σώμα πίεσης έως ονομαστικό βάθος 25, συμπεριλαμβανομένου του superflip σε απόσταση 20. Δεύτερον, με δίκαιη λογιστική, η αποφυγή του πλήρους ελέγχου του πίνακα ανά κλήση δίνει επιτάχυνση 218.855x, ενώ τα καθαρά οφέλη της λειτουργίας δέσμης και της παραμένουσας διεργασίας είναι 3.743x και 3.808x αντίστοιχα. Τρίτον, τα οφέλη αυτά αφορούν το κόστος φόρτωσης και ελέγχου του πίνακα και όχι την ίδια τη δύσκολη αναζήτηση, η οποία παραμένει το κυρίαρχο κόστος στις βαθιές μεμονωμένες περιπτώσεις. Οι εντολές παραγωγής των πινάκων και εκτέλεσης όλων των μετρήσεων της ενότητας παρατίθενται στο Παράρτημα~\ref{app:repro} και στο Παράρτημα~\ref{app:cli}.

\subsection{Δίκαιη λογιστική των επιταχύνσεων}
\label{sec:experiments-h48-fair}

Η διαδρομή H48 υποστηρίζει δύο τρόπους χρήσης του παραγόμενου πίνακα: την ελεγχόμενη λειτουργία (\eng{checked}), όπου κάθε κλήση επαληθεύει πλήρως τον πίνακα με \code{nissy_checkdata}, και την έμπιστη λειτουργία (\eng{trusted}), όπου ο πίνακας γίνεται δεκτός με βάση τα μεταδεδομένα παραγωγής του. Η έμπιστη λειτουργία στηρίζεται στα μεταδεδομένα του πίνακα \code{h48h7}· το ζήτημα της μη ντετερμινιστικής παραγωγής του πίνακα αυτού και η τεκμηριωμένη απόφαση διατήρησης των σχετικών τεκμηρίων συζητούνται στο Κεφάλαιο~\ref{chap:validation}, μαζί με την ανεξάρτητη διασταύρωση όλων των βαθιών ακριβών γραμμών.

Απαιτείται μια σημείωση μεθοδολογίας για τη λογιστική των μετρήσεων. Σε πρώτη καταγραφή χωρίς κοινή προθέρμανση, οι αποθηκευμένοι πολλαπλασιαστές ήταν 189.894x για τη σύγκριση ελεγχόμενης/έμπιστης λειτουργίας, 11.56x για τη δέσμη με ελεγχόμενο πίνακα, 38.198x για τη δέσμη με έμπιστο πίνακα και 40.16x για την παραμένουσα διεργασία. Οι τιμές αυτές χρέωναν το εφάπαξ ψυχρό κόστος εισαγωγής σελίδων (\eng{cold page-in}) του πίνακα των 3\,793\,842\,344 bytes εξ ολοκλήρου στο αργό σκέλος κάθε σύγκρισης. Τα αναθεωρημένα δίκαια τεκμήρια εκτελούν προθέρμανση πριν από κάθε χρονομετρημένο σκέλος και καταγράφουν σταθερή σειρά σκελών, ώστε κανένα σκέλος να μην πληρώνει το εφάπαξ κόστος μέσα στο μετρούμενο παράθυρο. Όλες οι τιμές που ακολουθούν προέρχονται από τα δίκαια τεκμήρια.

Στη δίκαιη σύγκριση ελεγχόμενης και έμπιστης λειτουργίας (Πίνακας~\ref{tab:h48-trusted-fair}), οι τέσσερις ελεγχόμενες και οι τέσσερις έμπιστες κλήσεις είναι όλες \code{exact} και επαληθευμένες: 42.058193~s για το ελεγχόμενο σκέλος έναντι 0.192174~s για το έμπιστο, δηλαδή επιτάχυνση 218.855x συνολικά και 232.606x σε σταθερή κατάσταση, εξαιρώντας την πρώτη χρονομετρημένη κλήση κάθε σκέλους. Η αναλογία είναι γνήσια και όχι τέχνημα κρυφής μνήμης, επειδή το ελεγχόμενο σκέλος επανασαρώνει ολόκληρο τον πίνακα \code{h48h7} σε κάθε κλήση. Η μέτρηση δείχνει ότι η καθυστέρηση βρίσκεται κυρίως στη σάρωση και επαλήθευση του μεγάλου πίνακα ανά διεργασία· δεν πρέπει να διαβαστεί ως επιτάχυνση αναζήτησης χειρότερης περίπτωσης για όλο τον χώρο καταστάσεων.

\begin{table}[htbp]\centering
{\small
\begin{tabular}{lrrr}
\hline
Mode & Total (s) & Steady-state (s) & Exact \\
\hline
Verified each call & 42.058193 & 31.739583 & 4 \\
Trusted generated table & 0.192174 & 0.136452 & 4 \\
Speedup & 218.855x & 232.606x & -- \\
\hline
\end{tabular}
}

\caption[Δίκαιη σύγκριση ελεγχόμενης και έμπιστης χρήσης του πίνακα h48h7]{Δίκαιη σύγκριση ελεγχόμενης και έμπιστης χρήσης του πίνακα \code{h48h7}, με κοινή προθέρμανση ώστε κανένα σκέλος να μη χρεώνεται το εφάπαξ ψυχρό κόστος φόρτωσης. Πηγή: \path{results/processed/h48_trusted_table_speedup_seed_2026_thesis_h48h7_fair.json}.}
\label{tab:h48-trusted-fair}
\end{table}

Η αντίστοιχη δίκαιη μέτρηση της λειτουργίας δέσμης με έμπιστο πίνακα (Πίνακας~\ref{tab:h48-batch-fair}) δείχνει ότι οι 12 επαναλήψεις της ίδιας ρηχής περίπτωσης παραμένουν \code{exact} και επαληθευμένες, με 0.67774~s σε ξεχωριστές έμπιστες κλήσεις έναντι 0.181076~s σε μία δέσμη, δηλαδή ρυθμαπόδοση (\eng{throughput}) 3.743x συνολικά και 3.112x σε σταθερή κατάσταση. Αφού αφαιρεθούν ο πλήρης έλεγχος του πίνακα και το εφάπαξ κόστος φόρτωσης, το πραγματικό κόστος διεργασίας και δέσμης για μικρές καταστάσεις παραμένει μετρήσιμο, αλλά είναι πολύ μικρότερο από τον αρχικά καταγεγραμμένο πολλαπλασιαστή.

\begin{table}[htbp]\centering
{\small
\begin{tabular}{lrrr}
\hline
Mode & Total (s) & Steady-state (s) & Exact \\
\hline
Single process per solve & 0.67774 & 0.516481 & 12 \\
Batch, table loaded once & 0.181076 & -- & 12 \\
Throughput speedup & 3.743x & 3.112x & -- \\
\hline
\end{tabular}
}

\caption[Δίκαιη μέτρηση του κόστους δέσμης με έμπιστο πίνακα]{Δίκαιη μέτρηση του κόστους δέσμης με έμπιστο πίνακα \code{h48h7} (κοινή προθέρμανση, 12 επαναλήψεις): ξεχωριστές διεργασίες έναντι μίας δέσμης με τον πίνακα φορτωμένο μία φορά. Πηγή: \path{results/processed/h48_batch_overhead_seed_2026_thesis_trusted_fair.json}.}
\label{tab:h48-batch-fair}
\end{table}

Τέλος, η δίκαιη σύγκριση ξεχωριστών κλήσεων και παραμένουσας διεργασίας (\eng{resident process}) στον Πίνακα~\ref{tab:h48-resident-fair} δίνει 0.325648~s για τις 8 ξεχωριστές κλήσεις έναντι 0.085528~s για τις 8 κλήσεις της παραμένουσας διεργασίας, δηλαδή 3.808x συνολικά και 6.55x σε σταθερή κατάσταση, εξαιρώντας την πρώτη κλήση κάθε σκέλους και την εκκίνηση της συνεδρίας. Η μέτρηση αφορά την επαναλαμβανόμενη χρήση και την αποφυγή της επανεκκίνησης διεργασίας και της επανεγκατάστασης του πίνακα για κάθε ρηχή κατάσταση.

\begin{table}[htbp]\centering
{\small
\begin{tabular}{lrrr}
\hline
Mode & Total (s) & Steady-state (s) & Exact \\
\hline
Separate native process & 0.325648 & 0.274602 & 8 \\
Resident native process & 0.085528 & 0.041924 & 8 \\
Speedup & 3.808x & 6.55x & -- \\
\hline
\end{tabular}
}

\caption[Δίκαιη σύγκριση ξεχωριστών και παραμένουσας εγγενούς διεργασίας]{Δίκαιη σύγκριση ξεχωριστών εγγενών διεργασιών και παραμένουσας εγγενούς διεργασίας για επαναλαμβανόμενες κλήσεις με έμπιστο πίνακα \code{h48h7} (κοινή προθέρμανση, 8 επαναλήψεις). Πηγή: \path{results/processed/h48_resident_oracle_seed_2026_thesis_h48h7_trusted_fair.json}.}
\label{tab:h48-resident-fair}
\end{table}

\subsection{Μετρήσεις πίεσης με τους πίνακες \code{h48h0} και \code{h48h7}}
\label{sec:experiments-h48-stress}

Η εκτέλεση πίεσης με τον μικρό πίνακα αναφοράς \code{h48h0} χρησιμοποιεί το ίδιο σταθερό σώμα περιπτώσεων με άμεση είσοδο κατάστασης. Ο Πίνακας~\ref{tab:h48h0-stress} είναι ξεχωριστός από το βασικό σώμα ακριβών τεκμηρίων, επειδή χρησιμοποιεί διαφορετική μηχανή επίλυσης (\eng{backend}) και προφίλ πίεσης. Η σήμανση παραμένει αυστηρή: οι γραμμές είναι \code{exact} επειδή η μηχανή H48 ολοκλήρωσε την αναζήτηση και οι λύσεις επαληθεύτηκαν ξανά από τον τοπικό επαληθευτή. Η ύπαρξη της διαδρομής αυτής δεν αλλάζει τα χρονικά όρια των μετρήσεων Python ούτε αποδεικνύει καθολική πρακτική κάλυψη κάθε δυνατής εισόδου 3×3×3.

\begin{table}[htbp]\centering
{\small
\begin{tabular}{lrrrrr}
\hline
Case & Depth & Initial LB & Status & Length & Seconds \\
\hline
solved & 0 & 0 & exact & 0 & 0.0 \\
shallow\_sequence & 3 & 3 & exact & 3 & 0.007653 \\
random\_1\_10 & 10 & 7 & exact & 9 & 0.015366 \\
random\_2\_15 & 15 & 9 & exact & 14 & 0.12254 \\
random\_3\_20 & 20 & 9 & exact & 17 & 12.721675 \\
\hline
\end{tabular}
}

\caption[Αποτελέσματα πίεσης της διαδρομής H48 με τον πίνακα h48h0]{Αποτελέσματα πίεσης της διαδρομής H48 με τον πίνακα αναφοράς \code{h48h0}: όλες οι γραμμές, συμπεριλαμβανομένης της ντετερμινιστικής περίπτωσης ονομαστικού βάθους 20, ολοκληρώνονται ως \code{exact}. Πηγή: \path{results/processed/optimal_3x3_seed_2026_stress_h48h0.json}.}
\label{tab:h48h0-stress}
\end{table}

Η ισχυρότερη εκτέλεση με τον πίνακα \code{h48h7} χρησιμοποιεί την ίδια άμεση είσοδο κατάστασης και προσθέτει δέκα ντετερμινιστικές περιπτώσεις ονομαστικού βάθους 25. Το αντίστοιχο αρχείο αποτελεσμάτων έχει 15/15 γραμμές \code{exact} και ανεξάρτητα επαληθευμένες, συμπεριλαμβανομένης της ντετερμινιστικής περίπτωσης βάθους 20 και των δέκα επιπλέον περιπτώσεων βάθους 25 (Πίνακας~\ref{tab:h48h7-stress}). Αυτό τεκμηριώνει ταχεία συμπεριφορά επιπέδου μαντείου στο αποθηκευμένο σώμα, όχι όμως ανεξάρτητη απόδειξη χρόνου χειρότερης περίπτωσης για κάθε δυνατή κατάσταση 3×3×3.

\begin{table}[htbp]\centering
{\small
\begin{tabular}{lrrrrr}
\hline
Case & Depth & Initial LB & Status & Length & Seconds \\
\hline
solved & 0 & 0 & exact & 0 & 0.0 \\
shallow\_sequence & 3 & 3 & exact & 3 & 0.01426 \\
random\_1\_10 & 10 & 7 & exact & 9 & 0.021346 \\
random\_2\_15 & 15 & 9 & exact & 14 & 0.123549 \\
random\_3\_20 & 20 & 9 & exact & 17 & 0.085337 \\
extra\_random\_1\_25 & 25 & 8 & exact & 18 & 1.400643 \\
extra\_random\_2\_25 & 25 & 8 & exact & 17 & 0.602382 \\
extra\_random\_3\_25 & 25 & 9 & exact & 17 & 0.349694 \\
extra\_random\_4\_25 & 25 & 9 & exact & 18 & 0.470633 \\
extra\_random\_5\_25 & 25 & 9 & exact & 18 & 0.731129 \\
extra\_random\_6\_25 & 25 & 9 & exact & 17 & 0.62516 \\
extra\_random\_7\_25 & 25 & 9 & exact & 18 & 1.729298 \\
extra\_random\_8\_25 & 25 & 9 & exact & 18 & 0.500993 \\
extra\_random\_9\_25 & 25 & 9 & exact & 17 & 0.522919 \\
extra\_random\_10\_25 & 25 & 9 & exact & 18 & 1.052717 \\
\hline
\end{tabular}
}

\caption[Αποτελέσματα πίεσης της διαδρομής H48 με τον πίνακα h48h7]{Αποτελέσματα πίεσης της διαδρομής H48 με τον πίνακα \code{h48h7}: 15/15 γραμμές \code{exact}, συμπεριλαμβανομένων δέκα ντετερμινιστικών περιπτώσεων ονομαστικού βάθους 25. Πηγή: \path{results/processed/optimal_3x3_seed_2026_stress_h48h7_oracle.json}.}
\label{tab:h48h7-stress}
\end{table}

Δοκιμάστηκε επίσης η παραγωγή του μεγαλύτερου πίνακα \code{h48h8} στο ίδιο μηχάνημα. Η προσπάθεια διακόπηκε μετά από εκτέλεση άνω των δύο ωρών χωρίς να παραχθεί το αρχείο \code{h48h8.bin}, με δείγματα εκτέλεσης να δείχνουν ότι ο χρόνος καταναλωνόταν μέσα στον εγγενή γεννήτορα H48 και όχι στο στρώμα Python. Επειδή δεν υπάρχει πλήρες αρχείο, άθροισμα ελέγχου ή μεταδεδομένα, η εργασία δεν χρησιμοποιεί το \code{h48h8} ως αποτέλεσμα και δεν αποδίδει σε αυτό κανέναν ισχυρισμό απόδοσης.

\subsection{Πιστοποίηση του δύσκολου σώματος}
\label{sec:experiments-h48-cert}

Η άμεση πιστοποίηση της διαδρομής μαντείου εκτελείται σε τέσσερις άμεσες περιπτώσεις \code{CubeState}. Η κρίσιμη γραμμή είναι το \code{superflip_distance_20}: όλες οι ακμές ανεστραμμένες, γωνίες και μεταθέσεις λυμένες, αναμενόμενη απόσταση 20 στη μετρική μισών στροφών (\eng{half-turn metric}, HTM). Το αποθηκευμένο αποτέλεσμα στον Πίνακα~\ref{tab:h48-cert} είναι \code{exact}, μήκους 20, επαληθευμένο, με μέγιστο χρόνο μεμονωμένης κλήσης 51.813928~s, κάτω από το καθορισμένο όριο των 90~s. Η ίδια εκτέλεση καταγράφει \code{backend_solve_seconds=44.591136} στο πεδίο σημειώσεων· η διαφορά είναι κυρίως κόστος διεργασίας, φόρτωσης και ελέγχου του πίνακα και όχι ένδειξη συντομότερης αναζήτησης.

\begin{table}[htbp]\centering
{\small
\begin{tabular}{lrrrr}
\hline
Case & Expected & Status & Length & Seconds \\
\hline
solved & 0 & exact & 0 & 0.0 \\
shallow\_r\_u\_f2 & 3 & exact & 3 & 6.146141 \\
deterministic\_depth\_25 & -- & exact & 18 & 7.091645 \\
superflip\_distance\_20 & 20 & exact & 20 & 51.813928 \\
\hline
\end{tabular}
}

\caption[Πιστοποίηση του μαντείου H48 σε ελεγχόμενη λειτουργία]{Πιστοποίηση της διαδρομής μαντείου H48 σε ελεγχόμενη λειτουργία (όριο 90~s): τέσσερις άμεσες περιπτώσεις, όλες \code{exact} και επαληθευμένες, με το superflip σε απόσταση 20. Πηγή: \path{results/processed/h48_oracle_certification_seed_2026_thesis.json}.}
\label{tab:h48-cert}
\end{table}

Η ίδια πιστοποίηση εκτελέστηκε και σε έμπιστη λειτουργία με προφόρτωση του πίνακα (Πίνακας~\ref{tab:h48-cert-preload}). Το αποτέλεσμα παραμένει 4/4 \code{exact} και επαληθευμένο, αλλά η γραμμή του superflip χρειάστηκε 158.440286~s. Πρόκειται για σημαντικό αρνητικό εύρημα για την ερμηνεία των επιταχύνσεων: η έμπιστη και προφορτωμένη διαδρομή αποφεύγει τον επαναλαμβανόμενο έλεγχο του πίνακα, αλλά δεν καθιστά τη δυσκολότερη αναζήτηση H48 ταχύτερη κατά τάξη μεγέθους. Επομένως οι ισχυρισμοί επιτάχυνσης δέκα φορών και άνω περιορίζονται ρητά σε περιπτώσεις κόστους φόρτωσης και ελέγχου πίνακα ή ρυθμαπόδοσης δέσμης.

\begin{table}[htbp]\centering
\resizebox{\textwidth}{!}{{\small
\begin{tabular}{lrrrr}
\hline
Case & Expected & Status & Length & Seconds \\
\hline
solved & 0 & exact & 0 & 0.0 \\
shallow\_r\_u\_f2 & 3 & exact & 3 & 8.960312 \\
deterministic\_depth\_25 & -- & exact & 18 & 56.078224 \\
superflip\_distance\_20 & 20 & exact & 20 & 158.440286 \\
\hline
\end{tabular}
}
}
\caption[Πιστοποίηση του μαντείου H48 σε έμπιστη λειτουργία με προφόρτωση]{Πιστοποίηση της διαδρομής μαντείου H48 σε έμπιστη λειτουργία με προφόρτωση (όριο 180~s): 4/4 \code{exact}, με τη γραμμή του superflip σε 158.440286~s. Πηγή: \path{results/processed/h48_oracle_certification_seed_2026_thesis_trusted_preload.json}.}
\label{tab:h48-cert-preload}
\end{table}

Καταγράφεται και τρίτη παραλλαγή της ίδιας πιστοποίησης, σε έμπιστη λειτουργία χωρίς προφόρτωση (όριο 300~s). Όπως δείχνει ο Πίνακας~\ref{tab:h48-cert-no-preload}, και εδώ οι τέσσερις γραμμές είναι \code{exact} και επαληθευμένες: η ντετερμινιστική περίπτωση ονομαστικού βάθους 25 λύνεται σε μήκος 18 σε 63.444134~s και το superflip σε μήκος 20 σε 91.629342~s. Η παραλλαγή αυτή ολοκληρώνει το δύσκολο σώμα χωρίς προφόρτωση και χρησιμοποιείται ως μέτρο αναφοράς στην ερμηνεία της Ενότητας~\ref{sec:experiments-interpretation}.

\begin{table}[htbp]\centering
\resizebox{\textwidth}{!}{{\small
\begin{tabular}{lrrrr}
\hline
Case & Expected & Status & Length & Seconds \\
\hline
solved & 0 & exact & 0 & 0.0 \\
shallow\_r\_u\_f2 & 3 & exact & 3 & 2.529118 \\
deterministic\_depth\_25 & -- & exact & 18 & 63.444134 \\
superflip\_distance\_20 & 20 & exact & 20 & 91.629342 \\
\hline
\end{tabular}
}
}
\caption[Πιστοποίηση του μαντείου H48 σε έμπιστη λειτουργία χωρίς προφόρτωση]{Πιστοποίηση της διαδρομής μαντείου H48 σε έμπιστη λειτουργία χωρίς προφόρτωση (όριο 300~s): 4/4 \code{exact}, με μέγιστο χρόνο 91.629342~s για το superflip. Πηγή: \path{results/processed/h48_oracle_certification_seed_2026_thesis_trusted_no_preload.json}.}
\label{tab:h48-cert-no-preload}
\end{table}

Τέλος, η διαδρομή της παραμένουσας διεργασίας δοκιμάστηκε όχι μόνο σε ρηχές εισόδους αλλά και στο ίδιο δύσκολο σώμα πιστοποίησης. Η εκτέλεση του Πίνακα~\ref{tab:h48-resident-cert} διατηρεί μία εγγενή διεργασία H48 και έναν απεικονισμένο στη μνήμη πίνακα \code{h48h7} για όλες τις γραμμές. Όλες οι γραμμές είναι \code{exact} και επαληθευμένες: η ντετερμινιστική περίπτωση βάθους 25 λύνεται σε μήκος 18 σε 72.071228~s και το superflip σε μήκος 20 σε 77.822389~s, με συνολικό χρόνο 154.755697~s. Αυτό είναι ισχυρότερο τεκμήριο για τη βελτιστοποιημένη διαδρομή μαντείου από τις ρηχές μικρομετρήσεις, αλλά παραμένει τεκμήριο σώματος και όχι εξαντλητική απόδειξη χρόνου χειρότερης περίπτωσης.

\begin{table}[htbp]\centering
\resizebox{\textwidth}{!}{{\small
\begin{tabular}{lrrrr}
\hline
Case & Expected & Status & Length & Seconds \\
\hline
solved & 0 & exact & 0 & 0.0 \\
shallow\_r\_u\_f2 & 3 & exact & 3 & 3.49045 \\
deterministic\_depth\_25 & -- & exact & 18 & 72.071228 \\
superflip\_distance\_20 & 20 & exact & 20 & 77.822389 \\
\hline
\end{tabular}
}
}
\caption[Πιστοποίηση του δύσκολου σώματος μέσω παραμένουσας διεργασίας]{Πιστοποίηση του δύσκολου σώματος μέσω παραμένουσας εγγενούς διεργασίας με έμπιστο πίνακα \code{h48h7} (όριο 240~s): 4/4 \code{exact}, με το superflip σε 77.822389~s. Πηγή: \path{results/processed/h48_resident_certification_seed_2026_thesis_h48h7_trusted.json}.}
\label{tab:h48-resident-cert}
\end{table}

\subsection{Λειτουργία δέσμης και δημόσιες διεπαφές του μαντείου}
\label{sec:experiments-h48-interfaces}

Για τη μέτρηση του κόστους ανά κλήση στην ελεγχόμενη λειτουργία, 12 επαναλήψεις της ίδιας ρηχής περίπτωσης εκτελέστηκαν ως ξεχωριστές διεργασίες και ως μία δέσμη με τον πίνακα φορτωμένο μία φορά (Πίνακας~\ref{tab:h48-batch-overhead}). Όλες οι γραμμές είναι \code{exact} και επαληθευμένες: 69.740540~s συνολικά για τις ξεχωριστές διεργασίες έναντι 6.033133~s για τη δέσμη. Στην ελεγχόμενη λειτουργία κάθε ξεχωριστή διεργασία πληρώνει την πλήρη φόρτωση και τον πλήρη έλεγχο του πίνακα, οπότε η διαφορά μετρά κυρίως αυτό το επαναλαμβανόμενο κόστος· η δίκαιη ποσοτικοποίηση του καθαρού κόστους δέσμης δίνεται στην Ενότητα~\ref{sec:experiments-h48-fair}.

\begin{table}[htbp]\centering
{\small
\begin{tabular}{lrr}
\hline
Mode & Total seconds & Exact rows \\
\hline
Single process per solve & 69.74054 & 12 \\
Batch table loaded once & 6.033133 & 12 \\
Throughput speedup & 11.56 & -- \\
\hline
\end{tabular}
}

\caption[Κόστος δέσμης σε ελεγχόμενη λειτουργία χωρίς κοινή προθέρμανση]{Κόστος δέσμης σε ελεγχόμενη λειτουργία, χωρίς κοινή προθέρμανση: 12 επαναλήψεις ως ξεχωριστές διεργασίες έναντι μίας δέσμης· η διαφορά κυριαρχείται από την επαναλαμβανόμενη φόρτωση και τον επαναλαμβανόμενο έλεγχο του πίνακα και υπερεκτιμά το καθαρό όφελος της δέσμης· η δίκαιη μέτρηση δίνεται στον Πίνακα~\ref{tab:h48-batch-fair}. Πηγή: \path{results/processed/h48_batch_overhead_seed_2026_thesis.json}.}
\label{tab:h48-batch-overhead}
\end{table}

Η ίδια διαδρομή δέσμης εκτίθεται ως δημόσια εντολή \code{rubik-optimal oracle}. Η εντολή δέχεται γραμμή προς γραμμή λυμένες καταστάσεις, ακολουθίες κινήσεων ή συμβολοσειρές ψηφίδων, μετατρέπει κάθε έγκυρη κατάσταση σε άμεση είσοδο H48 και καλεί την εγγενή μηχανή δέσμης, ώστε ο πίνακας \code{h48h7} να φορτωθεί μία φορά για όλο το σύνολο εισόδων. Το αποθηκευμένο τεκμήριο της διεπαφής γραμμής εντολών (Πίνακας~\ref{tab:h48-oracle-cli}) έχει τέσσερις εισόδους, όλες \code{exact} και ανεξάρτητα επαληθευμένες. Ο συνολικός χρόνος του περιβλήματος είναι 5.917467~s, ο χρόνος δέσμης που αναφέρει η ίδια η διεπαφή είναι 5.853855~s, και οι επιμέρους αναζητήσεις των μη λυμένων γραμμών απαιτούν εκατοστά του δευτερολέπτου. Η επαναλαμβανόμενη χρήση του μαντείου επωφελείται επομένως κατά τάξη μεγέθους, επειδή το ακριβό μέρος είναι η κοινή φόρτωση και ο έλεγχος του πίνακα.

\begin{table}[htbp]\centering
{\small
\begin{tabular}{llrr}
\hline
Case & Input & Distance & Seconds \\
\hline
solved & solved & 0 & 0.000000 \\
sequence\_shallow & sequence & 3 & 0.013599 \\
facelets\_shallow & facelets & 3 & 0.005634 \\
sequence\_mixed & sequence & 5 & 0.012297 \\
\hline
\end{tabular}
}

\caption[Τεκμήριο της εντολής δέσμης του μαντείου]{Τεκμήριο της δημόσιας εντολής \code{rubik-optimal oracle} σε λειτουργία δέσμης: τέσσερις είσοδοι (λυμένη κατάσταση, ακολουθίες, ψηφίδες), όλες \code{exact}, με τις επιμέρους αναζητήσεις σε εκατοστά του δευτερολέπτου. Πηγή: τεκμήριο του \code{scripts/run_h48_oracle_cli.py} για το προφίλ \code{thesis}.}
\label{tab:h48-oracle-cli}
\end{table}

Η αντίστοιχη εκτέλεση της ίδιας διεπαφής με έμπιστο πίνακα (Πίνακας~\ref{tab:h48-oracle-cli-trusted}) διατηρεί 4/4 γραμμές \code{exact} και επαληθευμένες και μειώνει τον συνολικό χρόνο του περιβλήματος σε 2.132448~s. Η μείωση είναι μικρότερη από τη συνθετική μέτρηση ανά κλήση, επειδή η διεπαφή χρησιμοποιεί ήδη φόρτωση δέσμης· παραμένει όμως χρήσιμη πρακτική βελτίωση για τη δημόσια χρήση της εντολής.

\begin{table}[htbp]\centering
{\small
\begin{tabular}{llrr}
\hline
Case & Input & Distance & Seconds \\
\hline
solved & solved & 0 & 0.000000 \\
sequence\_shallow & sequence & 3 & 1.120128 \\
facelets\_shallow & facelets & 3 & 0.009695 \\
sequence\_mixed & sequence & 5 & 0.884141 \\
\hline
\end{tabular}
}

\caption[Η εντολή του μαντείου με έμπιστο πίνακα]{Η εντολή \code{rubik-optimal oracle} με έμπιστο πίνακα \code{h48h7}: 4/4 γραμμές \code{exact} με μειωμένο συνολικό χρόνο περιβλήματος. Πηγή: τεκμήριο του \code{scripts/run_h48_oracle_cli.py} με επιλογή έμπιστου πίνακα.}
\label{tab:h48-oracle-cli-trusted}
\end{table}

Για επαναλαμβανόμενη χρήση χωρίς αναμονή ολόκληρης δέσμης, υποστηρίζεται λειτουργία ροής: η εντολή \code{rubik-optimal oracle --stream} διατηρεί μία παραμένουσα εγγενή διεργασία H48 και τυπώνει μία γραμμή JSON μόλις ολοκληρώνεται κάθε είσοδος. Το τεκμήριο του Πίνακα~\ref{tab:h48-oracle-stream} έχει 4/4 γραμμές \code{exact} και επαληθευμένες, με λυμένη κατάσταση, ακολουθίες κινήσεων και είσοδο ψηφίδων. Η πρώτη μη λυμένη γραμμή απορροφά το κόστος ετοιμότητας της παραμένουσας διεργασίας, ενώ οι επόμενες χρησιμοποιούν τον ήδη απεικονισμένο πίνακα.

\begin{table}[htbp]\centering
{\small
\begin{tabular}{llrr}
\hline
Case & Input & Distance & Seconds \\
\hline
solved & solved & 0 & 0.000000 \\
sequence\_shallow & sequence & 3 & 1.640139 \\
facelets\_shallow & facelets & 3 & 0.012159 \\
sequence\_mixed & sequence & 5 & 2.509894 \\
\hline
\end{tabular}
}

\caption[Λειτουργία ροής του μαντείου με παραμένουσα διεργασία]{Λειτουργία ροής της εντολής \code{rubik-optimal oracle} με παραμένουσα εγγενή διεργασία και έμπιστο πίνακα: 4/4 γραμμές \code{exact}, με την πρώτη μη λυμένη γραμμή να απορροφά το κόστος ετοιμότητας. Πηγή: τεκμήριο του \code{scripts/run_h48_oracle_stream.py}.}
\label{tab:h48-oracle-stream}
\end{table}

Η ίδια υλοποίηση εκτίθεται και ως προγραμματιστική διεπαφή σε επίπεδο πακέτου, μέσω της κλάσης \code{FastOptimalOracle}, ώστε ο όρος «μαντείο» να μην αφορά μόνο περιβλήματα γραμμής εντολών. Το τεκμήριο του Πίνακα~\ref{tab:fast-oracle-api} περιλαμβάνει λυμένη κατάσταση, ακολουθία κινήσεων, είσοδο ψηφίδων και ντετερμινιστική κατάσταση ονομαστικού βάθους 10· οι 4/4 γραμμές είναι \code{exact} και επαληθευμένες, με μέγιστο χρόνο 2.664585~s. Το ίδιο αρχείο συνδέει και το τεκμήριο της πιστοποίησης του δύσκολου σώματος, ώστε η προγραμματιστική απόδειξη να μην αποκόπτεται από τις δυσκολότερες γραμμές. Και εδώ ο χρόνος εκτέλεσης παραμένει εμπειρικό τεκμήριο και όχι εξαντλητική χρονική απόδειξη για κάθε πιθανή κατάσταση.

\begin{table}[htbp]\centering
{\small
\begin{tabular}{llrrr}
\hline
Case & Input & Expected & Length & Seconds \\
\hline
solved & cube\_state & 0 & 0 & 0.0 \\
sequence\_shallow & sequence & 3 & 3 & 2.393554 \\
facelets\_shallow & facelets & 3 & 3 & 0.020547 \\
deterministic\_depth\_10 & deterministic\_scramble & -- & 10 & 2.664585 \\
\hline
\end{tabular}
}

\caption[Τεκμήριο της προγραμματιστικής διεπαφής του μαντείου]{Τεκμήριο της προγραμματιστικής διεπαφής \code{FastOptimalOracle} με έμπιστο πίνακα \code{h48h7}: 4/4 γραμμές \code{exact} με μέγιστο χρόνο 2.664585~s. Πηγή: τεκμήριο του \code{scripts/run_fast_optimal_oracle_api.py}.}
\label{tab:fast-oracle-api}
\end{table}

\subsection{Συμβόλαιο ισχυρισμών}
\label{sec:experiments-h48-contract}

Για να μην παραμείνει το όριο των ισχυρισμών μόνο ως αφηγηματικό σχόλιο, παράγεται χωριστό τεκμήριο συμβολαίου που ελέγχει αυτόματα τις κρίσιμες ιδιότητες της διαδρομής μαντείου. Οι ουσιώδεις έλεγχοι είναι πέντε. Πρώτον, η τεκμηρίωση και οι επικεφαλίδες του ενσωματωμένου nissy-core επιβεβαιώνουν ότι η διαδρομή H48 είναι βέλτιστος επιλυτής στη μετρική μισών στροφών και ότι το ψευδώνυμο \code{optimal} δείχνει στο \code{h48h7}. Δεύτερον, το περίβλημα C καλεί τη \code{nissy_solve} με σημαία κανονικής λειτουργίας και μέγιστο βάθος 20. Τρίτον, το περίβλημα Python και η κλάση \code{FastOptimalOracle} εκτελούν άμεση μετατροπή κατάστασης κυβιδίων (\eng{cubies}), έλεγχο φυσικής εγκυρότητας, χρήση παραμένουσας μηχανής \code{h48h7} με μεταδεδομένα παραγωγής και ανεξάρτητη επαλήθευση κάθε λύσης. Τέταρτον, το συμβόλαιο συνδέει τα αποθηκευμένα εμπειρικά τεκμήρια της ενότητας. Πέμπτον, το πεδίο της καθολικής απόδειξης ταχείας εκτέλεσης για κάθε κατάσταση παραμένει σκόπιμα αρνητικό: η διεπαφή δέχεται κάθε φυσικά έγκυρη κατάσταση 3×3×3 και τηρεί σύμβαση αποκλειστικά ακριβών αποτελεσμάτων για όσες γραμμές ολοκληρώνονται, χωρίς να μετατρέπει τους χρόνους του σώματος σε θεώρημα χρόνου χειρότερης περίπτωσης. Ο πλήρης πίνακας των ελέγχων του συμβολαίου παρατίθεται στο Παράρτημα~\ref{app:repro}.

\subsection{Εξωτερική διαδρομή Nissy}
\label{sec:experiments-h48-nissy}

Συμπληρωματικά προς τη διαδρομή H48, εγκαταστάθηκε το δημόσιο σύνολο πινάκων \code{pt_nxopt31_HTM} του Nissy~2.x~\cite{trontoNissyCoreH48}, με λήψη μόνο της αντίστοιχης καταχώρισης ZIP και επαλήθευση του αθροίσματος ελέγχου CRC. Δεν πρόκειται για εναλλακτικό πίνακα H48, αλλά για χωριστή εξωτερική διαδρομή \code{nissy-optimal}. Το σώμα πίεσης της διαδρομής αυτής (Πίνακας~\ref{tab:nissy-optimal-stress}) έχει 5/5 γραμμές \code{exact} και επαληθευμένες· η ντετερμινιστική περίπτωση ονομαστικού βάθους 20 αναγνωρίζεται σε απόσταση 17 και λύνεται σε 23.734412~s με 2 νήματα, υπό φορτίο γενικής χρήσης.

Η ίδια διαδρομή δεν αντικαθιστά την πιστοποίηση των δύσκολων περιπτώσεων: στις αποθηκευμένες οριοθετημένες δοκιμές, η εντολή \code{solve optimal} του εξωτερικού \code{nissy} δεν ολοκλήρωσε το superflip εντός 120~s με 2 νήματα ούτε εντός 240~s με 8 νήματα (τεκμήρια στο \path{results/processed/}). Η κατάσταση \code{timeout} στις δοκιμές αυτές καταγράφει μόνο τη μη ολοκλήρωση εντός του ορίου χρόνου στο συγκεκριμένο μηχάνημα· δεν αποτελεί φράγμα απόστασης. Το βέλτιστο σύνολο πινάκων του Nissy προστίθεται έτσι ως δεύτερη βελτιστοποιημένη ακριβής μηχανή και ως τεκμήριο για το σώμα πίεσης, ενώ η πιστοποίηση μέσω παραμένουσας διεργασίας H48 παραμένει το ισχυρότερο αποθηκευμένο \emph{ενσωματωμένο} τεκμήριο για το superflip σε απόσταση 20. Η ισχυρότερη \emph{εγγενής} απόδειξη του ίδιου αποτελέσματος είναι η εξαντλητική απόδειξη της Ενότητας~\ref{sec:experiments-superflip}.

\begin{table}[htbp]\centering
{\small
\begin{tabular}{lrrrrr}
\hline
Case & Depth & Initial LB & Status & Length & Seconds \\
\hline
solved & 0 & 0 & exact & 0 & 5.1e-05 \\
shallow\_sequence & 3 & 3 & exact & 3 & 2.083463 \\
random\_1\_10 & 10 & 7 & exact & 9 & 2.107859 \\
random\_2\_15 & 15 & 9 & exact & 14 & 1.747235 \\
random\_3\_20 & 20 & 9 & exact & 17 & 23.734412 \\
\hline
\end{tabular}
}

\caption[Σώμα πίεσης της εξωτερικής διαδρομής Nissy]{Σώμα πίεσης της εξωτερικής διαδρομής \code{nissy-optimal} με το δημόσιο σύνολο πινάκων \code{pt_nxopt31_HTM}: 5/5 γραμμές \code{exact}, με 2 νήματα. Πηγή: τεκμήριο του \code{scripts/run_3x3_optimal.py} με μηχανή \code{nissy-optimal}.}
\label{tab:nissy-optimal-stress}
\end{table}

\section{Εγγενής απόδειξη βελτιστότητας του superflip}
\label{sec:experiments-superflip}

Πέρα από τις πιστοποιήσεις μέσω H48 και Nissy, η εργασία περιλαμβάνει πλήρη εγγενή απόδειξη ότι η απόσταση του superflip είναι ακριβώς 20, χωρίς καμία εξάρτηση από H48 ή Nissy. Η απόδειξη συνδυάζει δύο σκέλη: την ήδη επαληθευμένη εγγενή λύση μήκους 20 από τη διαδρομή δύο φάσεων Kociemba, η οποία δίνει το άνω φράγμα, και μια εξαντλητική απόδειξη κάτω φράγματος. Η μέθοδος του δεύτερου σκέλους --- συμμετρικά ανηγμένος πίνακας φάσης~1 \eng{FlipUDSlice} 16 συμμετριών με 64\,430 κλάσεις ισοδυναμίας κατά Kociemba~\cite{kociembaImplementationDetails}, παραδεκτό φράγμα τριών αξόνων κατά Reid~\cite{reid1997optimal} και αποκοπή ρίζας από τον σταθεροποιητή συμμετρίας 48~στοιχείων του superflip --- περιγράφεται αναλυτικά στο Κεφάλαιο~\ref{chap:implementation}.

Η αποθηκευμένη εκτέλεση είναι η εξής: ένας εξαντλητικός IDA* ενός φράγματος (\code{--mode optimal-ida}) που αναζητά απευθείας τη λυμένη κατάσταση εξάντλησε το όριο μήκους 19 χωρίς να φθάσει σε λυμένο φύλλο. Αυτό αποδεικνύει ότι δεν υπάρχει λύση μήκους 19 ή μικρότερου· μαζί με τη λύση μήκους 20 αποδεικνύεται εγγενώς απόσταση ακριβώς 20. Η παραδεκτότητα της ευρετικής επικυρώθηκε με σύγκριση \eng{byte-for-byte} έναντι του ωμού πίνακα BFS φάσης~1 σε όλες τις 959\,761\,462 γνωστές εγγραφές με απόσταση φάσης~1 έως 9, με μηδέν αναντιστοιχίες (τεκμήριο: \path{results/processed/native_sym_phase1_verify_seed_2026_thesis.json}). Επιπλέον, η ίδια αναζήτηση βρίσκει σε ρηχές περιπτώσεις βέλτιστη λύση στη σωστή απόσταση και αποκλείει κάθε λύση ένα βήμα χαμηλότερα. Το πλήθος των κόμβων είναι ντετερμινιστικό και ανεξάρτητο από το πλήθος νημάτων και την κατάτμηση εργασιών, οπότε το αποτέλεσμα είναι αναπαραγώγιμο. Ο μόνος περιορισμός είναι ο χρόνος και όχι η ορθότητα: η εξαντλητική απόδειξη για αυτή τη χειρότερη περίπτωση απαιτεί περίπου 87 λεπτά στο μηχάνημα των 8 πυρήνων και 16~GiB. Ο Πίνακας~\ref{tab:superflip-proof} συνοψίζει τα αποθηκευμένα τεκμήρια, με όλες τις τιμές να προέρχονται από τα αποθηκευμένα αρχεία αποτελεσμάτων.

\begin{table}[htbp]\centering
% Generated by scripts/generate_superflip_proof_table.py -- do not edit by hand.
{\small
\begin{tabular}{L{0.34\textwidth}L{0.58\textwidth}}
\hline
Quantity & Value \\
\hline
Solver / mode & \code{kociemba_reid_optimal_ida} (\code{--mode optimal-ida}) \\
Heuristic & FlipUDSlice 16-symmetry phase-1 (max 12) + three-axis $\max(p_{ud},p_{rl},p_{fb})$ \\
Root symmetry mask & U, U2 (48-element stabilizer of superflip) \\
Uses H48/Nissy & false \\
Target bound (length) & 19 \\
Status & \code{lower_bound} \\
Solved leaf found & false \\
Timed out / node-limited & false / false \\
Proves no solution $\le$ 19 & true \\
Expanded nodes & 11\,203\,278\,121 \\
Generated nodes & 149\,020\,204\,380 \\
Runtime (8 threads) & 5200.8 s \\
Companion bound-18 proof & proves no solution $\le$ 18 (true), 248.5 s \\
FlipUDSlice symmetry table & 64\,430 classes $\times$ 2\,187 twist $=$ 140\,908\,410 phase-1 entries (exact to depth 12) \\
\hline
\end{tabular}
}

\caption[Τεκμήρια της εγγενούς απόδειξης κάτω φράγματος του superflip]{Αποθηκευμένα τεκμήρια της εγγενούς εξαντλητικής απόδειξης κάτω φράγματος για το superflip: παραμετροποίηση της αναζήτησης, ντετερμινιστικά πλήθη κόμβων και αποτέλεσμα αποκλεισμού κάθε λύσης μήκους έως 19. Πηγή: τεκμήρια της εκτέλεσης \code{--mode optimal-ida}.}
\label{tab:superflip-proof}
\end{table}

\section{Χαρτοφυλάκιο ακριβών επιλυτών}
\label{sec:experiments-portfolio}

Η πρακτική συνέπεια της ύπαρξης δύο ακριβών μηχανών είναι ότι δεν χρειάζεται να επιλεγεί ένας μοναδικός επιλυτής για όλες τις περιπτώσεις. Η κλάση \code{PortfolioOptimalOracle} επαναχρησιμοποιεί επανεπικυρωμένα ακριβή πιστοποιητικά για ήδη αποδεδειγμένες καταστάσεις, δοκιμάζει πρώτα τη διαδρομή \code{nissy-optimal} με περιορισμένο χρόνο και, αν αυτή δεν επιστρέψει ακριβές επαληθευμένο αποτέλεσμα, μεταβιβάζει την ίδια κατάσταση στο παραμένον μαντείο H48. Η διαδρομή αυτή δεν αλλάζει τον κανόνα ακρίβειας: η κατάσταση \code{exact} αποδίδεται μόνο όταν η επιλεγμένη μηχανή επιστρέψει λύση που επαληθεύεται ανεξάρτητα, ή όταν ένα αποθηκευμένο ακριβές και επαληθευμένο πιστοποιητικό επανελεγχθεί πάνω στην ίδια κατάσταση.

Σε εκτέλεση χαμηλής προτεραιότητας υπό φορτίο γενικής χρήσης, το σώμα με προτεραιότητα στο Nissy χρησιμοποιεί μία διεργασία δέσμης Nissy, έχει 5/5 γραμμές \code{exact} και μέγιστο χρόνο δέσμης 10.097333~s (Πίνακας~\ref{tab:portfolio-nissy-first}). Η στοχευμένη γραμμή εφεδρείας του superflip λύνει το superflip απόστασης 20 μέσω παραμένουσας διεργασίας H48 σε 135.679618~s με 4 νήματα H48 (Πίνακας~\ref{tab:portfolio-superflip-fallback}). Μετά την αποθήκευση του αντίστοιχου ακριβούς πιστοποιητικού, η ίδια κατάσταση επιστρέφεται από την κρυφή μνήμη πιστοποιητικών σε 0.012885~s χωρίς νέα ακριβή αναζήτηση (Πίνακας~\ref{tab:portfolio-superflip-cache}). Το αποτέλεσμα είναι τεκμήριο μηχανικής σύνθεσης χαρτοφυλακίου, όχι απόδειξη ότι κάθε πιθανή κατάσταση έχει πρακτικό χρόνο χειρότερης περίπτωσης.

\begin{table}[htbp]\centering
{\small
\begin{tabular}{lrrrr}
\hline
Case & Source depth & Status & Length & Seconds \\
\hline
solved & 0 & exact & 0 & 0.0 \\
shallow\_sequence & 3 & exact & 3 & 10.097333 \\
random\_1\_10 & 10 & exact & 9 & 10.097333 \\
random\_2\_15 & 15 & exact & 14 & 10.097333 \\
random\_3\_20 & 20 & exact & 17 & 10.097333 \\
\hline
\end{tabular}
}

\caption[Σώμα του χαρτοφυλακίου με προτεραιότητα στο Nissy]{Σώμα του χαρτοφυλακίου ακριβών επιλυτών με προτεραιότητα στο Nissy, σε εκτέλεση χαμηλής προτεραιότητας: 5/5 γραμμές \code{exact} μέσω μίας διεργασίας δέσμης. Πηγή: τεκμήριο \code{lowload} του \code{PortfolioOptimalOracle}.}
\label{tab:portfolio-nissy-first}
\end{table}

\begin{table}[htbp]\centering
{\small
\begin{tabular}{lrrrr}
\hline
Case & Source depth & Status & Length & Seconds \\
\hline
superflip\_distance\_20 & 0 & exact & 20 & 135.679618 \\
\hline
\end{tabular}
}

\caption[Γραμμή εφεδρείας του χαρτοφυλακίου για το superflip]{Γραμμή εφεδρείας του χαρτοφυλακίου για το superflip: ακριβής λύση μήκους 20 μέσω παραμένουσας διεργασίας H48 σε 135.679618~s. Πηγή: τεκμήριο \code{lowload} του \code{PortfolioOptimalOracle}.}
\label{tab:portfolio-superflip-fallback}
\end{table}

\begin{table}[htbp]\centering
{\small
\begin{tabular}{lrrrr}
\hline
Case & Source depth & Status & Length & Seconds \\
\hline
superflip\_distance\_20 & 0 & exact & 20 & 0.012885 \\
\hline
\end{tabular}
}

\caption[Επιστροφή του superflip από την κρυφή μνήμη πιστοποιητικών]{Επιστροφή του superflip από την κρυφή μνήμη πιστοποιητικών του χαρτοφυλακίου σε 0.012885~s, μετά την επανεπικύρωση του αποθηκευμένου ακριβούς πιστοποιητικού. Πηγή: τεκμήριο \code{lowload} του \code{PortfolioOptimalOracle}.}
\label{tab:portfolio-superflip-cache}
\end{table}

\section{Σύνοψη επιλυτών}
\label{sec:experiments-summary}

Ο Πίνακας~\ref{tab:benchmark-summary} συνοψίζει το κύριο σώμα μετρήσεων ανά επιλυτή και δείχνει ότι οι επιλυτές έχουν διαφορετικούς ρόλους. Ο οριοθετημένος Korf/IDA* σε Python επιστρέφει \code{exact} μόνο στις περιπτώσεις όπου ολοκληρώνει την αναζήτηση: με τον γωνιακό πίνακα προτύπων και όριο βάθους 10 αποδεικνύει 11 από τις 13 περιπτώσεις, συμπεριλαμβανομένης της τυχαίας περίπτωσης ονομαστικού βάθους 10, όπου βρίσκει ακριβή απόσταση 9. Ο εγγενής βέλτιστος επιλυτής με γωνιακούς και ακμικούς πίνακες προτύπων αποτελεί χωριστή διαδρομή τεκμηρίων με μεγαλύτερο βάθος· στο αποθηκευμένο σώμα αποδεικνύει ακριβείς λύσεις για τη λυμένη κατάσταση, τη ρηχή ακολουθία και τις τυχαίες περιπτώσεις βάθους 10 και 15 (Πίνακας~\ref{tab:optimal-status}). Η διαδρομή H48 είναι τρίτη, πιο εξειδικευμένη διαδρομή ακριβούς αναζήτησης: λύνει ακριβώς και τη ντετερμινιστική περίπτωση βάθους 20, τα 15/15 του σώματος πίεσης και το superflip (Πίνακες~\ref{tab:h48h0-stress}, \ref{tab:h48h7-stress} και~\ref{tab:h48-cert}).

Ο εξωτερικός προσαρμογέας Kociemba~\cite{muodovKociembaPkg} έχει πλήρη κάλυψη επαληθευμένων λύσεων στο κύριο σώμα, αλλά όλες οι λύσεις του παραμένουν \code{non_exact}. Η εγγενής υλοποίηση Kociemba χρησιμοποιεί φασικούς πίνακες και επαληθεύει όλες τις 13 περιπτώσεις με κατάσταση \code{non_exact}. Η εγγενής διαδρομή Thistlethwaite εκτελεί τις τέσσερις μειώσεις υποομάδων ως άπληστη κάθοδο πάνω στους ακριβείς πίνακες αποστάσεων συμπλόκων (\eng{coset}), τερματίζει εγγυημένα χωρίς όρια χρόνου αναζήτησης και επαληθεύει επίσης όλες τις 13 περιπτώσεις. Οι επιτυχίες αυτές δεν μετατρέπουν τις λύσεις σε αποδείξεις βελτιστότητας· δείχνουν πρακτική τετραφασική επίλυση στο σταθερό σώμα μετρήσεων.

\begin{table}[htbp]\centering
{\small
\begin{tabular}{L{0.28\textwidth}rrrL{0.16\textwidth}}
\hline
Solver & Cases & Verified & Median runtime (s) & Statuses \\
\hline
Korf & 13 & 11 & 0.030163 & exact, lower\_bound \\
Koc. native & 13 & 13 & 0.003210 & non\_exact \\
Koc. adapter & 13 & 13 & 0.000473 & non\_exact \\
Thistle. & 13 & 13 & 0.013516 & non\_exact \\
\hline
\end{tabular}
}

\caption[Σύνοψη του κύριου σώματος μετρήσεων ανά επιλυτή]{Σύνοψη του κύριου σώματος μετρήσεων ανά επιλυτή: πλήθος περιπτώσεων, επαληθευμένες λύσεις, διάμεσος χρόνος εκτέλεσης και εμφανιζόμενες καταστάσεις αποτελέσματος. Πηγή: \path{results/processed/summary_seed_2026_thesis.json}.}
\label{tab:benchmark-summary}
\end{table}

Ο Πίνακας~\ref{tab:algorithm-status} συνοψίζει συμπληρωματικά την κατάσταση υλοποίησης και το είδος εγγύησης κάθε αλγοριθμικής διαδρομής της εργασίας.

\begin{table}[htbp]\centering
{\small
\begin{tabular}{L{0.17\textwidth}L{0.29\textwidth}L{0.25\textwidth}L{0.17\textwidth}}
\hline
Algorithm & Implementation status & Optimality status & Evidence \\
\hline
BFS & Shallow exhaustive search & exact within configured depth & tests/results \\
Korf/IDA* & Scoped Python IDA* with projection tables and native corner/edge PDB lower bounds & exact when completed & tests + PDB \\
Korf native optimal & C++ full-cube IDA* with complete corner and 6-edge PDBs & exact when completed & optimal evidence \\
Kociemba native & Scoped two-phase: exact sym-reduced phase-1 heuristic plus phase-2 projection pruning & non-exact verified when completed & tests + tables \\
Kociemba adapter & Optional external two-phase adapter & non-exact verified solution & verifier/results \\
Thistlethwaite native & Four-phase greedy descent over exact BFS stage-distance tables & non-exact verified, guaranteed termination & tests + tables \\
\hline
\end{tabular}
}

\caption[Κατάσταση υλοποίησης και εγγυήσεις αλγοριθμικών διαδρομών]{Κατάσταση υλοποίησης, εγγύηση βελτιστότητας και πηγή τεκμηρίων για κάθε αλγοριθμική διαδρομή της εργασίας. Πηγή: παραγόμενος πίνακας κατάστασης αλγορίθμων του προφίλ \code{thesis}.}
\label{tab:algorithm-status}
\end{table}

\section{Συγκεντρωτική εικόνα αποτελεσμάτων}
\label{sec:experiments-overview}

Τα παρακάτω σχήματα και πίνακες δίνουν διαφορετικές όψεις του ίδιου αρχείου αποτελεσμάτων: κατανομή καταστάσεων, χρόνο εκτέλεσης, μήκος επαληθευμένων λύσεων και κόμβους αναζήτησης όπου ο επιλυτής εκθέτει τέτοια μέτρηση.

Το Σχήμα~\ref{fig:status-counts} και ο Πίνακας~\ref{tab:solver-status-counts} δείχνουν την κατανομή των καταστάσεων αποτελέσματος ανά επιλυτή. Η κατανομή αυτή είναι πιο σημαντική από το απλό πλήθος επαληθευμένων λύσεων. Ένας επιλυτής που δίνει επαληθευμένες \code{non_exact} λύσεις σε όλες τις περιπτώσεις είναι πρακτικά χρήσιμος, αλλά δεν απαντά στο ερώτημα της ακριβούς απόστασης. Ένας επιλυτής που δίνει \code{exact} λύσεις σε λιγότερες περιπτώσεις παρέχει ισχυρότερη εγγύηση εκεί όπου επιτυγχάνει. Το σώμα μετρήσεων επομένως δεν κατατάσσει τους επιλυτές με μονοδιάστατη βαθμολογία.

\begin{figure}[htbp]\centering
% Κατανομή καταστάσεων αποτελέσματος ανά επιλυτή (ομαδοποιημένο ραβδόγραμμα).
% Περιεχόμενο σχήματος· το περιβάλλον figure και η λεζάντα ορίζονται στο
% κεφάλαιο 9. Δεδομένα: figures/status_counts_data.tex. Μηδενικές ράβδοι
% παραλείπονται ανά σειρά.
\begin{tikzpicture}
\begin{axis}[
  width=0.85\textwidth,
  height=0.55\textwidth,
  ybar,
  bar width=11pt,
  ylabel={Πλήθος περιπτώσεων},
  symbolic x coords={Korf, Koc. native, Koc. adapter, Thistle.},
  xtick={Korf, Koc. native, Koc. adapter, Thistle.},
  ymin=0, ymax=16,
  ytick={0,2,4,6,8,10,12,14,16},
  enlarge x limits=0.18,
  ymajorgrids,
  grid style={dotted, gray!50},
  legend pos=north west,
  legend cell align=left,
  legend style={font=\footnotesize},
]
\addplot[fill=blue!60, draw=blue!50!black] coordinates {
  (Korf,11)
};
\addlegendentry{\texttt{exact}}
\addplot[fill=orange!75, draw=orange!60!black] coordinates {
  (Korf,2)
};
\addlegendentry{\texttt{lower\_bound}}
\addplot[fill=gray!55, draw=gray!70!black] coordinates {
  (Koc. native,13)
  (Koc. adapter,13)
  (Thistle.,13)
};
\addlegendentry{\texttt{non\_exact}}
\end{axis}
\end{tikzpicture}

\caption[Κατανομή καταστάσεων αποτελέσματος ανά επιλυτή]{Κατανομή καταστάσεων αποτελέσματος ανά επιλυτή στο κύριο σώμα μετρήσεων (\code{exact}, \code{lower_bound}, \code{non_exact}).}
\label{fig:status-counts}
\end{figure}

\begin{table}[htbp]\centering
\begin{tabular}{p{0.24\textwidth}rrrrrr}
\hline
Solver & Exact & Non-exact & Lower bound & Timeout & N/A & Failed \\
\hline
Korf & 11 & 0 & 2 & 0 & 0 & 0 \\
Koc. native & 0 & 13 & 0 & 0 & 0 & 0 \\
Koc. adapter & 0 & 13 & 0 & 0 & 0 & 0 \\
Thistle. & 0 & 13 & 0 & 0 & 0 & 0 \\
\hline
\end{tabular}

\caption[Πλήθη καταστάσεων αποτελέσματος ανά επιλυτή]{Πλήθη καταστάσεων αποτελέσματος ανά επιλυτή στο κύριο σώμα μετρήσεων. Πηγή: \path{results/raw/benchmarks_seed_2026_thesis.jsonl}.}
\label{tab:solver-status-counts}
\end{table}

Το Σχήμα~\ref{fig:runtime-depth} δείχνει την κλιμάκωση του χρόνου εκτέλεσης του οριοθετημένου επιλυτή Korf στις ρηχές περιπτώσεις, ενώ το Σχήμα~\ref{fig:runtime-by-solver} και ο Πίνακας~\ref{tab:solver-runtime} συνοψίζουν τους χρόνους όλων των επιλυτών. Οι χρόνοι των οριοθετημένων επιλυτών περιλαμβάνουν τη λειτουργία της υλοποίησης σε Python και την πρόσβαση στους αποθηκευμένους πίνακες αποκοπής (\eng{pruning tables})· δεν πρέπει να συγκρίνονται απευθείας με βελτιστοποιημένες εγγενείς υλοποιήσεις της βιβλιογραφίας.

\begin{figure}[htbp]\centering
% Χρόνος εκτέλεσης του επιλυτή Korf ως προς το βάθος ανακατέματος (1-8).
% Περιεχόμενο σχήματος· το περιβάλλον figure και η λεζάντα ορίζονται στο
% κεφάλαιο 9. Δεδομένα: figures/runtime_depth_data.tex.
\begin{tikzpicture}
\begin{axis}[
  width=0.85\textwidth,
  height=0.55\textwidth,
  xlabel={Βάθος ανακατέματος},
  ylabel={Χρόνος εκτέλεσης (s)},
  xmin=0.5, xmax=8.5,
  xtick={1,2,3,4,5,6,7,8},
  ymin=0, ymax=0.07,
  scaled y ticks=false,
  yticklabel style={/pgf/number format/fixed, /pgf/number format/precision=3},
  grid=major,
  grid style={dotted, gray!50},
  legend pos=north west,
  legend cell align=left,
  legend style={font=\footnotesize},
]
\addplot[blue, thick, mark=*, mark size=1.8pt] coordinates {
  (1,0.001663)
  (2,0.002298)
  (3,0.017847)
  (4,0.024893)
  (5,0.060181)
  (6,0.022632)
  (7,0.030163)
  (8,0.059619)
};
\addlegendentry{Korf (IDA*)}
\end{axis}
\end{tikzpicture}

\caption[Χρόνος εκτέλεσης του επιλυτή Korf ως προς το βάθος διατάραξης]{Χρόνος εκτέλεσης του επιλυτή Korf ως προς το βάθος διατάραξης (1--8) στο κύριο σώμα μετρήσεων με σπόρο 2026.}
\label{fig:runtime-depth}
\end{figure}

\begin{figure}[htbp]\centering
% Διάμεσος και μέγιστος χρόνος εκτέλεσης ανά επιλυτή (λογαριθμική κλίμακα).
% Περιεχόμενο σχήματος· το περιβάλλον figure και η λεζάντα ορίζονται στο
% κεφάλαιο 9. Δεδομένα: figures/runtime_by_solver_data.tex. Οι ελάχιστοι
% χρόνοι δεν απεικονίζονται (μηδενική τιμή σε λογαριθμικό άξονα)· δίνονται
% στον συνοδευτικό πίνακα.
\begin{tikzpicture}
\begin{axis}[
  width=0.85\textwidth,
  height=0.55\textwidth,
  ybar,
  bar width=13pt,
  ymode=log,
  log origin=infty,
  ylabel={Χρόνος (s), λογαριθμική κλίμακα},
  symbolic x coords={Korf, Koc. native, Koc. adapter, Thistle.},
  xtick={Korf, Koc. native, Koc. adapter, Thistle.},
  ymin=1e-4, ymax=100,
  enlarge x limits=0.18,
  ymajorgrids,
  grid style={dotted, gray!50},
  legend pos=north west,
  legend cell align=left,
  legend style={font=\footnotesize},
]
\addplot[fill=blue!60, draw=blue!50!black] coordinates {
  (Korf,0.030163)
  (Koc. native,0.003210)
  (Koc. adapter,0.000473)
  (Thistle.,0.013516)
};
\addlegendentry{Διάμεσος}
\addplot[fill=red!65, draw=red!50!black] coordinates {
  (Korf,2.023518)
  (Koc. native,10.638481)
  (Koc. adapter,0.097952)
  (Thistle.,1.514541)
};
\addlegendentry{Μέγιστος}
\end{axis}
\end{tikzpicture}

\caption[Διάμεσος και μέγιστος χρόνος εκτέλεσης ανά επιλυτή]{Διάμεσος και μέγιστος χρόνος εκτέλεσης ανά επιλυτή στο κύριο σώμα μετρήσεων (λογαριθμική κλίμακα)· οι ελάχιστοι χρόνοι αναφέρονται στον Πίνακα~\ref{tab:solver-runtime}.}
\label{fig:runtime-by-solver}
\end{figure}

\begin{table}[htbp]\centering
\input{tables/solver_runtime_summary.tex}
\caption[Χρόνοι εκτέλεσης ανά επιλυτή]{Ελάχιστος, διάμεσος και μέγιστος χρόνος εκτέλεσης ανά επιλυτή στο κύριο σώμα μετρήσεων. Πηγή: \path{results/raw/benchmarks_seed_2026_thesis.jsonl}.}
\label{tab:solver-runtime}
\end{table}

Το Σχήμα~\ref{fig:solution-length-depth} και ο Πίνακας~\ref{tab:solver-length} συνοψίζουν τα μήκη των επαληθευμένων λύσεων. Τα μήκη είναι πιο σταθερή πληροφορία από τους χρόνους, αλλά χρειάζονται και αυτά το πλαίσιο της κατάστασης αποτελέσματος: μια επαληθευμένη λύση μήκους 8 με κατάσταση \code{non_exact} δεν είναι καλύτερη ή χειρότερη από μια ακριβή λύση χωρίς γνώση της αρχικής κατάστασης και του αποδεικτικού πλαισίου --- απλώς συνοδεύεται από διαφορετική εγγύηση.

\begin{figure}[htbp]\centering
% Διάμεσο μήκος λύσης ανά βάθος ανακατέματος και επιλυτή. Περιεχόμενο
% σχήματος· το περιβάλλον figure και η λεζάντα ορίζονται στο κεφάλαιο 9.
% Δεδομένα: figures/solution_length_depth_data.tex. Η σειρά Korf
% τερματίζει στο βάθος 10 (βαθύτερες περιπτώσεις χωρίς γραμμή διαμέσου).
\begin{tikzpicture}
\begin{axis}[
  width=0.85\textwidth,
  height=0.62\textwidth,
  xlabel={Βάθος ανακατέματος},
  ylabel={Διάμεσο μήκος λύσης},
  xmin=-0.5, xmax=20.5,
  xtick={0,5,10,15,20},
  minor x tick num=4,
  ymin=0, ymax=40,
  ytick={0,5,10,15,20,25,30,35,40},
  grid=major,
  grid style={dotted, gray!50},
  legend pos=north west,
  legend cell align=left,
  legend style={font=\footnotesize},
]
% Διαγώνιος αναφοράς y = x (βέλτιστο μήκος ίσο με το ονομαστικό βάθος).
\addplot[forget plot, gray, thick] coordinates {(0,0) (20,20)}
  node[pos=0.85, sloped, below, font=\scriptsize, text=gray] {βέλτιστο};
\addplot[blue, thick, mark=*, mark size=1.6pt] coordinates {
  (0,0) (1,1) (2,2) (3,3) (4,4) (5,5) (6,4) (7,7) (8,8) (10,18) (15,22) (20,23)
};
\addlegendentry{Kociemba (native)}
\addplot[red, thick, mark=square*, mark size=1.6pt] coordinates {
  (0,0) (1,1) (2,2) (3,3) (4,4) (5,5) (6,4) (7,7) (8,8) (10,9) (15,21) (20,21)
};
\addlegendentry{Kociemba (adapter)}
\addplot[green!55!black, thick, mark=triangle*, mark size=2pt] coordinates {
  (0,0) (1,1) (2,2) (3,3) (4,4) (5,5) (6,4) (7,7) (8,8) (10,9)
};
\addlegendentry{Korf}
\addplot[orange!90!black, thick, mark=diamond*, mark size=2pt] coordinates {
  (0,0) (1,1) (2,2) (3,3) (4,5) (5,17) (6,4) (7,25) (8,32) (10,31) (15,32) (20,26)
};
\addlegendentry{Thistlethwaite}
\end{axis}
\end{tikzpicture}

\caption[Διάμεσο μήκος λύσης ανά βάθος διατάραξης και επιλυτή]{Διάμεσο μήκος λύσης ανά βάθος διατάραξης και επιλυτή· οι βέλτιστοι επιλυτές ακολουθούν κατά κανόνα τη διαγώνιο του ονομαστικού βάθους έως το βάθος στο οποίο ολοκληρώνουν, ενώ ο Thistlethwaite επιστρέφει επαληθευμένες αλλά μη βέλτιστες λύσεις.}
\label{fig:solution-length-depth}
\end{figure}

\begin{table}[htbp]\centering
\begin{tabular}{p{0.28\textwidth}rrrr}
\hline
Solver & Verified & Min & Median & Max \\
\hline
Korf & 11 & 0 & 4.000000 & 9 \\
Koc. native & 13 & 0 & 5.000000 & 23 \\
Koc. adapter & 13 & 0 & 5.000000 & 21 \\
Thistle. & 13 & 0 & 5.000000 & 32 \\
\hline
\end{tabular}

\caption[Μήκη επαληθευμένων λύσεων ανά επιλυτή]{Πλήθος επαληθευμένων λύσεων και ελάχιστο, διάμεσο και μέγιστο μήκος λύσης ανά επιλυτή. Πηγή: \path{results/raw/benchmarks_seed_2026_thesis.jsonl}.}
\label{tab:solver-length}
\end{table}

Οι κόμβοι αναζήτησης (Πίνακας~\ref{tab:solver-nodes}) είναι διαθέσιμοι για τους εγγενείς επιλυτές αναζήτησης. Δεν είναι διαθέσιμοι από τον εξωτερικό προσαρμογέα, επειδή το πακέτο δεν εκθέτει αυτή την εσωτερική μέτρηση μέσα από τη διεπαφή που χρησιμοποιείται εδώ. Η κενή τιμή δεν είναι μηδέν· σημαίνει μη διαθέσιμη μέτρηση. Η διάκριση αυτή αποτρέπει τη λανθασμένη ανάγνωση ότι ο προσαρμογέας έλυσε χωρίς αναζήτηση.

\begin{table}[htbp]\centering
\begin{tabular}{p{0.32\textwidth}rr}
\hline
Solver & Median expanded & Median generated \\
\hline
Korf & 5 & 78 \\
Koc. native & 101 & 1527 \\
Koc. adapter & -- & -- \\
Thistle. & 5 & 5 \\
\hline
\end{tabular}

\caption[Κόμβοι αναζήτησης ανά επιλυτή]{Διάμεσοι κόμβοι επέκτασης και παραγωγής ανά επιλυτή· η κενή τιμή του προσαρμογέα σημαίνει μη διαθέσιμη μέτρηση. Πηγή: \path{results/raw/benchmarks_seed_2026_thesis.jsonl}.}
\label{tab:solver-nodes}
\end{table}

\section{Μήτρα περιπτώσεων και μη ολοκληρώσεις}
\label{sec:experiments-matrix}

Η μήτρα του Πίνακα~\ref{tab:status-matrix} δείχνει ανά περίπτωση την κατάσταση κάθε επιλυτή και, όπου υπάρχει λύση, το μήκος της. Οι συντομογραφίες είναι: E για \code{exact}, N για \code{non_exact} και LB για \code{lower_bound}, δηλαδή αποδεδειγμένο κάτω φράγμα χωρίς λύση. Η αποθηκευμένη εκτέλεση δεν παρήγαγε καμία γραμμή \code{timeout}: οι βαθύτερες μη ολοκληρώσεις του Korf καταγράφονται ως \code{lower_bound}, επειδή προκύπτουν από εξάντληση του ορίου βάθους με τεκμηριωμένο φράγμα και όχι από εξάντληση χρονικού ορίου χωρίς αποτέλεσμα.

Η μήτρα είναι ο πιο άμεσος τρόπος να εντοπισθεί αν οι μη ολοκληρώσεις συγκεντρώνονται σε συγκεκριμένες κατηγορίες. Οι ρηχές περιπτώσεις λειτουργούν ως έλεγχος ότι οι επιλυτές δεν αποτυγχάνουν σε απλές εισόδους, ενώ οι βαθύτερες τυχαίες περιπτώσεις αποκαλύπτουν πού σταματούν οι οριοθετημένες αναζητήσεις. Η μήτρα βοηθά επίσης να μη διαβάζεται απομονωμένα ο διάμεσος χρόνος εκτέλεσης: ένας επιλυτής μπορεί να εμφανίζει χαμηλό διάμεσο χρόνο επειδή επιστρέφει γρήγορα στις εύκολες περιπτώσεις, ενώ εξακολουθεί να μην ολοκληρώνει στις δύσκολες.

\begin{table}[htbp]\centering
\begin{tabular}{p{0.23\textwidth}rllll}
\hline
Case & Depth & Korf & Koc. native & Koc. adapter & Thistle. \\
\hline
solved & 0 & E/0 & N/0 & N/0 & N/0 \\
shallow\_1 & 1 & E/1 & N/1 & N/1 & N/1 \\
shallow\_2 & 2 & E/2 & N/2 & N/2 & N/2 \\
shallow\_3 & 3 & E/3 & N/3 & N/3 & N/3 \\
shallow\_4 & 4 & E/4 & N/4 & N/4 & N/5 \\
shallow\_5 & 5 & E/5 & N/5 & N/5 & N/29 \\
shallow\_6 & 6 & E/4 & N/4 & N/4 & N/4 \\
shallow\_7 & 7 & E/7 & N/7 & N/7 & N/25 \\
shallow\_8 & 8 & E/8 & N/8 & N/8 & N/32 \\
random\_0\_5 & 5 & E/5 & N/5 & N/5 & N/5 \\
random\_1\_10 & 10 & E/9 & N/18 & N/9 & N/31 \\
random\_2\_15 & 15 & LB & N/22 & N/21 & N/32 \\
random\_3\_20 & 20 & LB & N/23 & N/21 & N/26 \\
\hline
\end{tabular}

\caption[Μήτρα κατάστασης ανά περίπτωση και επιλυτή]{Μήτρα κατάστασης ανά περίπτωση και επιλυτή στο κύριο σώμα μετρήσεων (E~= \code{exact}, N~= \code{non_exact}, LB~= \code{lower_bound})· όπου υπάρχει λύση αναγράφεται και το μήκος της. Πηγή: \path{results/raw/benchmarks_seed_2026_thesis.jsonl}.}
\label{tab:status-matrix}
\end{table}

Οι περιπτώσεις \code{timeout} καταγράφονται χωριστά, επειδή δεν πρέπει να μετατρέπονται σε σιωπηρή επιτυχία ούτε να παρουσιάζονται ως ακριβής απόσταση. Όπως τεκμηριώνει ρητά ο Πίνακας~\ref{tab:timeout-cases}, η συγκεκριμένη εκτέλεση του κύριου σώματος δεν παρήγαγε καμία τέτοια γραμμή.

\begin{table}[htbp]\centering
{\small
\begin{tabular}{L{0.16\textwidth}L{0.18\textwidth}rrL{0.30\textwidth}}
\hline
Solver & Case & Depth & Expanded & Note \\
\hline
-- & -- & -- & -- & No timeout rows \\
\hline
\end{tabular}
}

\caption[Γραμμές εξάντλησης χρονικού ορίου του κύριου σώματος μετρήσεων]{Γραμμές \code{timeout} του κύριου σώματος μετρήσεων: η αποθηκευμένη εκτέλεση δεν παρήγαγε καμία, καθώς οι βαθύτερες μη ολοκληρώσεις καταγράφονται ως \code{lower_bound}. Πηγή: \path{results/raw/benchmarks_seed_2026_thesis.jsonl}.}
\label{tab:timeout-cases}
\end{table}

\section{Πλήρης κατανομή του κύβου 2×2×2}
\label{sec:experiments-pocket}

Για τον κύβο 2×2×2 εκτελέστηκε πλήρης εξαντλητική BFS στο κανονικοποιημένο μοντέλο της μετρικής μισών στροφών, η οποία επισκέφθηκε και τις 3\,674\,160 καταστάσεις και βρήκε μέγιστη απόσταση 11. Η εντολή παραγωγής παρατίθεται στο Παράρτημα~\ref{app:repro}.

Ο Πίνακας~\ref{tab:pocket-distribution} δείχνει πόσες καταστάσεις βρίσκονται σε κάθε απόσταση από τη λυμένη κατάσταση μέσα στο κανονικοποιημένο μοντέλο. Η κατανομή κορυφώνεται στην απόσταση 9, με 1\,887\,748 καταστάσεις, ενώ στη μέγιστη απόσταση 11 βρίσκονται μόλις 2644 καταστάσεις. Η εικόνα αυτή είναι αναμενόμενη σε μεγάλους χώρους καταστάσεων: οι πολύ ρηχές καταστάσεις είναι μόνο όσες παράγονται από μικρές ακολουθίες κινήσεων, ενώ η μεγάλη μάζα βρίσκεται κοντά στο μέγιστο βάθος. Το ουσιώδες για την εργασία δεν είναι η σύγκριση της κατανομής αυτής με τον πλήρη κύβο 3×3×3, αλλά η επίδειξη ότι η διαδικασία πλήρους ακριβούς κατανομής ολοκληρώνεται και επαληθεύεται.

\begin{table}[htbp]\centering
\begin{tabular}{rr}
\hline
Distance & States \\
\hline
0 & 1 \\
1 & 9 \\
2 & 54 \\
3 & 321 \\
4 & 1847 \\
5 & 9992 \\
6 & 50136 \\
7 & 227536 \\
8 & 870072 \\
9 & 1887748 \\
10 & 623800 \\
11 & 2644 \\
\hline
\end{tabular}

\caption[Πλήρης κατανομή αποστάσεων του κύβου 2×2×2]{Πλήρης κατανομή αποστάσεων του κύβου 2×2×2 στο κανονικοποιημένο μοντέλο της μετρικής μισών στροφών: πλήθος καταστάσεων ανά απόσταση από τη λυμένη κατάσταση, με κορύφωση στην απόσταση 9 και μέγιστη απόσταση 11. Πηγή: \path{results/processed/pocket_cube_summary_seed_2026_thesis.json}.}
\label{tab:pocket-distribution}
\end{table}

Οι αντιπροσωπευτικές λύσεις του Πίνακα~\ref{tab:pocket-representatives} επιλέγονται για να δείξουν ότι η πλήρης κατανομή δεν είναι μόνο συγκεντρωτική καταμέτρηση. Για συγκεκριμένες καταστάσεις, το σύστημα επιστρέφει λύση, μήκος, κόμβους επέκτασης και παραγωγής, κατάσταση \code{exact} και ένδειξη επαλήθευσης. Οι περιπτώσεις αυτές είναι λίγες, αλλά συνδέουν την κατανομή αποστάσεων με πραγματικές ακολουθίες κινήσεων.

\begin{table}[htbp]\centering
\begin{tabular}{p{0.2\textwidth}rp{0.34\textwidth}p{0.16\textwidth}}
\hline
Case & Length & Solution & Status \\
\hline
pocket\_rep\_1 & 1 & R' & exact \\
pocket\_rep\_2 & 2 & U' R' & exact \\
pocket\_rep\_3 & 3 & F' U' R' & exact \\
pocket\_rep\_4 & 4 & F2 R U' R' & exact \\
pocket\_rep\_5 & 6 & F U R U' R' F' & exact \\
\hline
\end{tabular}

\caption[Αντιπροσωπευτικές λύσεις του κύβου 2×2×2]{Αντιπροσωπευτικές ακριβείς λύσεις του κύβου 2×2×2 με μήκος, ακολουθία λύσης και κατάσταση αποτελέσματος. Πηγή: \path{results/processed/pocket_cube_summary_seed_2026_thesis.json}.}
\label{tab:pocket-representatives}
\end{table}

\section{Ερμηνεία}
\label{sec:experiments-interpretation}

Στις ρηχές περιπτώσεις, ο οριοθετημένος επιλυτής Korf/IDA* επιστρέφει ακριβείς λύσεις που συμφωνούν με την BFS όπου υπάρχει εξαντλητικό σημείο αναφοράς. Με τον γωνιακό πίνακα προτύπων, ο ίδιος επιλυτής αποδεικνύει και την τυχαία περίπτωση ονομαστικού βάθους 10 μέσα στο όριο βάθους 10. Με τους ακμικούς πίνακες προτύπων και τον εγγενή βέλτιστο επιλυτή, η χωριστή διαδρομή ακριβών τεκμηρίων αποδεικνύει επιπλέον την τυχαία περίπτωση ονομαστικού βάθους 15 ως απόσταση 14 (Πίνακας~\ref{tab:optimal-status}). Με τη διαδρομή H48, το σώμα πίεσης αποδεικνύει τη ντετερμινιστική περίπτωση βάθους 20 ως ακριβή λύση μήκους 17 (Πίνακας~\ref{tab:h48h0-stress}), η πιστοποίηση με τον πίνακα \code{h48h7} αποδεικνύει το superflip ως ακριβή λύση μήκους 20 (Πίνακας~\ref{tab:h48-cert}), η έμπιστη πιστοποίηση χωρίς προφόρτωση ολοκληρώνει το δύσκολο σώμα με μέγιστο χρόνο 91.629342~s (Πίνακας~\ref{tab:h48-cert-no-preload}), και η λειτουργία δέσμης αφαιρεί το επαναλαμβανόμενο κόστος φόρτωσης του πίνακα στις επαναληπτικές κλήσεις του μαντείου. Πρόκειται για ουσιαστική πρόοδο στον κύβο 3×3×3, η οποία όμως δεν αποτελεί από μόνη της απόδειξη ότι κάθε αυθαίρετη κατάσταση θα ολοκληρωθεί τοπικά σε συγκεκριμένο πρακτικό χρόνο. Υπενθυμίζεται επίσης ότι τα τεκμήρια που βασίζονται στον πίνακα \code{h48h7} συνοδεύονται από την επιφύλαξη μη ντετερμινιστικής παραγωγής του πίνακα και από την ανεξάρτητη διασταύρωση των βαθιών ακριβών γραμμών, όπως τεκμηριώνεται στο Κεφάλαιο~\ref{chap:validation}.

Η εγγενής διαδρομή Kociemba χρησιμοποιεί αναζήτηση στον χώρο συντεταγμένων και στις δύο φάσεις και επιστρέφει επαληθευμένες λύσεις σε όλες τις περιπτώσεις του κύριου σώματος. Η εγγενής διαδρομή Thistlethwaite ολοκληρώνει επίσης όλες τις περιπτώσεις, χωρίς καθόλου οριοθετημένη αναζήτηση: κάθε φάση είναι κάθοδος καθοδηγούμενη από ακριβή πίνακα αποστάσεων. Ο προσαρμογέας Kociemba παραμένει πρακτικό σημείο αναφοράς και καταγράφεται επίσης ως \code{non_exact}. Η πλήρης μελέτη του κύβου 2×2×2 λειτουργεί ως μικρό ολοκληρωμένο παράδειγμα εξαντλητικής βέλτιστης αναζήτησης, όχι ως υποκατάστατο των διαδρομών Korf, PDB και H48 του 3×3×3.

Η βασική ερμηνεία είναι ότι το κόστος της απόδειξης διαφέρει από το κόστος της εύρεσης μιας λύσης. Ο προσαρμογέας βρίσκει γρήγορα λύσεις επειδή χρησιμοποιεί ώριμη πρακτική μέθοδο. Η οριοθετημένη IDA* χρειάζεται να αποκλείσει όλες τις συντομότερες λύσεις για να αποδώσει \code{exact}. Η εγγενής Kociemba και η εγγενής Thistlethwaite βρίσκονται σε ενδιάμεση περιοχή: έχουν σχεδίαση διαδοχικών σταδίων και λύνουν το σταθερό σώμα μετρήσεων --- η Thistlethwaite μάλιστα τερματίζει εγγυημένα σε κάθε έγκυρη κατάσταση χάρη στους ακριβείς φασικούς πίνακες --- αλλά οι λύσεις τους δεν συνοδεύονται από απόδειξη ελάχιστου μήκους. Επομένως δεν υποστηρίζουν ισχυρισμό καθολικής βέλτιστης επίλυσης του 3×3×3 ούτε πλήρους αναπαραγωγής των ιστορικών πινάκων ελιγμών.

Τα αποτελέσματα δείχνουν επίσης γιατί η εργασία δεν βασίζεται σε έναν μόνο επιλυτή. Αν υπήρχε μόνο ο προσαρμογέας, το σύστημα θα έλυνε πρακτικά τις περιπτώσεις, αλλά δεν θα επιδείκνυε ακριβή αναζήτηση. Αν υπήρχε μόνο η IDA*, η εργασία θα είχε καθαρότερη εγγύηση αλλά λιγότερες λυμένες βαθιές περιπτώσεις. Ο συνδυασμός επιτρέπει να παρουσιασθεί η διαφορά ανάμεσα στις καταστάσεις \code{exact}, \code{non_exact} και \code{lower_bound} με πραγματικά καταγεγραμμένες γραμμές: ο Kociemba και ο Thistlethwaite λύνουν πρακτικά το σώμα μετρήσεων, ενώ η διαδρομή Korf/IDA* δείχνει πού ολοκληρώνεται και πού όχι η βέλτιστη απόδειξη.

Τέλος, κάθε αριθμός που εμφανίζεται στους πίνακες του κεφαλαίου προέρχεται από αποθηκευμένο αρχείο αποτελεσμάτων ή παραγόμενα μεταδεδομένα, και οι πίνακες παράγονται αυτόματα από τα ίδια αρχεία. Έτσι το κείμενο παραμένει ελέγξιμο: κάθε ισχυρισμός του κεφαλαίου μπορεί να αναχθεί στο αντίστοιχο τεκμήριο με τις διαδικασίες του Παραρτήματος~\ref{app:repro}.
